\chapter{Finite Fields at Play}

\begin{center}\itshape
A mathematician, like a painter or a poet, is a maker of patterns. If his patterns are more permanent than theirs, it is
because they are made with ideas.\\
\upshape --- G.~H.~Hardy, \emph{A Mathematician's Apology} (1940)
\end{center}

In 1974 Marsha Jean Falco, a population geneticist, was studying epilepsy in German shepherd dogs. To look for patterns in the
genetic data she drew little symbols on cards, each card showing a number of shapes with a colour and a shading. The cards
turned out to make a good puzzle, and in 1991 they were published as the game \emph{SET}. There are $81$ cards, one for each choice of four attributes
(number, colour, shape, shading) with three values each. Three cards form a \emph{set} when, for every attribute, they are all
the same or all different. The players lay out twelve cards and race to spot a set.

Write the three values of each attribute as $0,1,2$, so that a card is a vector of $\mathbb F_3^4$. Three numbers of $\mathbb F_3$
are all equal or all different exactly when their sum is $0$, so three cards $x,y,z$ form a set exactly when $x+y+z=0$: they
are the three points of a line of the affine space $\mathbb F_3^4$. A collection of cards without a set is a \emph{cap}. The
largest cap has $20$ cards (Pellegrino, 1971), which is why a table of $21$ cards always holds a set.

The \emph{cap set problem} asks the same question in $\mathbb F_3^n$. For decades the best upper bound was a little smaller than
$3^n$: Meshulam proved $O(3^n/n)$ in 1995 with Fourier analysis, and Bateman and Katz, after years of work, shaved off a tiny bit
more in 2012. Then in May 2016 Ernie Croot, Vsevolod Lev and Péter Pál Pach posted a paper on the group $(\Z/4\Z)^n$ with a new
polynomial trick, and about a week later Jordan Ellenberg and Dion Gijswijt, independently, adapted it to $\mathbb F_3^n$. Their proof
was three pages long and gave $2.756^n$: an exponentially small fraction of the space. Terence Tao rewrote it the same week on
his blog in the language of ``slice rank'' that we use in Problem~\ref{pb:capset}.

This chapter is about the polynomial method and its cousins. Finite fields are small enough to count everything and rigid
enough that a polynomial of low degree cannot vanish in too many places. Out of these two facts come point counts on curves,
additive combinatorics, Kakeya sets, and objects that look random without being random.

\begin{think}
Lay out SET cards one by one, never completing a set. How long can you go? Try it in $\mathbb F_3^2$ (nine cards with two
attributes) and in $\mathbb F_3^3$. Do you believe that, in $\mathbb F_3^n$, a cap can fill a fixed proportion of the space,
say a hundredth of it, for every $n$?
\end{think}

\section{Counting points}

Over a finite field every equation has a number of solutions, and the problems of this section compute some of them. In most
of them $p$ is an odd prime, $\chi(a)=\bigl(\frac ap\bigr)$ is the Legendre symbol with $\chi(0)=0$, and $e(t)=e^{2i\pi t/p}$.

\begin{exercise}[$\star$]\label{pb:conics}
\emph{Conics.} Let $p$ be an odd prime.
\begin{enumerate}[label=(\alph*)]
\item Let $a\ne0$ and $b,c$ be in $\mathbb F_p$, and $D=b^2-4ac$. Show that
\[
\sum_{x\in\mathbb F_p}\chi(ax^2+bx+c)=\begin{cases}-\chi(a)&\text{if }D\ne0,\\(p-1)\chi(a)&\text{if }D=0.\end{cases}
\]
\item Let $Q(x,y)=ax^2+bxy+cy^2$ with $D=b^2-4ac\ne0$. Show that the equation $Q(x,y)=1$ has exactly $p-\chi(D)$ solutions in
$\mathbb F_p^2$.
\item Deduce that a binary quadratic form with non-zero discriminant takes every non-zero value of $\mathbb F_p$, and that
$x^2+y^2=1$ has $p-\chi(-1)$ solutions.
\end{enumerate}
\end{exercise}

\begin{exercise}[$\star\star$]\label{pb:sumsquares}
\emph{Sums of squares over $\mathbb F_p$.} Let $p$ be an odd prime, $n\ge1$, and $N_n(c)$ the number of $x\in\mathbb F_p^n$ with
$x_1^2+\dots+x_n^2=c$. Write $\eta=\chi$ for short.
\begin{enumerate}[label=(\alph*)]
\item Show that for $c\ne0$,
\[
N_n(c)=\begin{cases}p^{n-1}-p^{n/2-1}\,\eta\bigl((-1)^{n/2}\bigr)&n\text{ even},\\[2pt]
p^{n-1}+p^{(n-1)/2}\,\eta\bigl((-1)^{(n-1)/2}c\bigr)&n\text{ odd},\end{cases}
\]
and find $N_n(0)$.
\item Draw $x$ uniformly in $\mathbb F_p^n$. Show that the total variation distance between the law of $x_1^2+\dots+x_n^2$ and
the uniform law on $\mathbb F_p$ is
\[
\Bigl(1-\frac1p\Bigr)p^{-n/2}\quad(n\text{ even}),\qquad\frac{p-1}2\,p^{-(n+1)/2}\quad(n\text{ odd}).
\]
\end{enumerate}
\end{exercise}

\begin{exercise}[$\star\star\star$]\label{pb:fermatcubic}
\emph{Gauss's cubic.} Let $p\equiv1\pmod3$ be prime and $N$ the number of $(x,y)\in\mathbb F_p^2$ with $x^3+y^3=1$.
\begin{enumerate}[label=(\alph*)]
\item Let $\chi$ be a character of order $3$ of $\mathbb F_p^*$, extended by $\chi(0)=0$, and for characters $\alpha,\beta$
let $J(\alpha,\beta)=\sum_{u+v=1}\alpha(u)\beta(v)$ (with the convention that the trivial character takes the value $1$ at $0$).
Show that $N=p-2+2\operatorname{Re}J(\chi,\chi)$.
\item Let $g(\alpha)=\sum_{t\in\mathbb F_p}\alpha(t)e(t)$. Show that $g(\alpha)g(\beta)=J(\alpha,\beta)g(\alpha\beta)$ when
$\alpha\beta$ is not trivial, that $|g(\chi)|^2=p$, and that $g(\chi)^3=p\,J(\chi,\chi)$.
\item Write $\omega=e^{2i\pi/3}$, so that $J(\chi,\chi)=a'+b'\omega$ with $a',b'\in\Z$. Show that $a'\equiv-1$ and
$b'\equiv0\pmod 3$.
\item Conclude: $N=p-2+a$, where $a$ is the integer with $a\equiv1\pmod3$ and $4p=a^2+27b^2$ for some integer $b$. Check $p=7$
and $p=13$, and show that $|N-(p-2)|\le2\sqrt p$.
\end{enumerate}
\end{exercise}

\begin{figure}[ht]
\centering
\fig{ch05-cubic}
\caption{For every prime $p\equiv1\pmod3$ below $3000$, the error $a=N-(p-2)$ in the number of points of $x^3+y^3=1$, divided by
$2\sqrt p$ (Problem~\ref{pb:fermatcubic}). It never leaves $[-1,1]$, and it does not settle anywhere: the points fill the band
with a density that peaks near the edges.}
\label{fig:cubic}
\end{figure}

\begin{exercise}[$\star\star\star$]\label{pb:nilpotent}
\emph{Counting nilpotent matrices.} Let $N_n$ be the number of nilpotent matrices in $M_n(\mathbb F_q)$, with $N_0=1$.
\begin{enumerate}[label=(\alph*)]
\item (Fitting.) For $A\in M_n(\mathbb F_q)$ acting on $V=\mathbb F_q^n$, show that $V=\ker A^n\oplus\operatorname{im}A^n$,
that both are $A$-stable, that $A$ is nilpotent on the first and invertible on the second, and that this is the only
decomposition of $V$ into two $A$-stable subspaces with these properties.
\item Deduce that $\displaystyle\sum_{k=0}^n\frac{N_k}{|GL_k(\mathbb F_q)|}=\frac{q^{n^2}}{|GL_n(\mathbb F_q)|}$ for every $n\ge0$.
\item (Fine and Herstein, 1958.) Prove that $N_n=q^{n^2-n}$. So a random matrix is nilpotent with probability exactly $q^{-n}$.
\item (Steinberg.) How many unipotent matrices ($A-I$ nilpotent) does $GL_n(\mathbb F_q)$ contain?
\end{enumerate}
\end{exercise}

\begin{exercise}[$\star\star$]\label{pb:rankdist}
\emph{The rank of a random matrix.}
\begin{enumerate}[label=(\alph*)]
\item Show that the number of matrices of rank $r$ in $M_n(\mathbb F_q)$ is
\[
\prod_{i=0}^{r-1}\frac{(q^n-q^i)^2}{q^r-q^i}.
\]
\item Let $A$ be uniform in $M_n(\mathbb F_q)$. Show that for each $k\ge0$, as $n\to\infty$,
\[
P(\operatorname{corank}A=k)\longrightarrow q^{-k^2}\,\frac{\prod_{j>k}(1-q^{-j})}{\prod_{j=1}^k(1-q^{-j})}.
\]
\item Deduce Durfee's identity: for $|x|<1$, writing $(x)_k=(1-x)(1-x^2)\cdots(1-x^k)$,
\[
\sum_{k\ge0}\frac{x^{k^2}}{(x)_k^{\,2}}=\prod_{j\ge1}\frac1{1-x^j},
\]
first for $x=\frac1q$, then for all $x$.
\end{enumerate}
\end{exercise}

\begin{exercise}[$\star$]\label{pb:flags}
\emph{Flags and the field with one element.} Let $V=\mathbb F_q^n$ and $[m]_q=1+q+\dots+q^{m-1}$, $[n]_q!=[1]_q[2]_q\cdots[n]_q$.
A \emph{complete flag} is a chain $0=V_0\subset V_1\subset\dots\subset V_n=V$ with $\dim V_i=i$.
\begin{enumerate}[label=(\alph*)]
\item Show that $V$ has $[n]_q!$ complete flags.
\item Let $U$ be a subspace of dimension $k$. Show that $U$ has exactly $q^{k(n-k)}$ complements $W$ (subspaces with $U\oplus W=V$).
\item Show that $|GL_n(\mathbb F_q)|=q^{\binom n2}(q-1)^n\,[n]_q!$, by letting $GL_n$ act on flags.
\item Set $q=1$ in the formulas of (a) and (b). What do flags and complements become, and what does (c) say about the
symmetric group?
\end{enumerate}
\end{exercise}

\section{The polynomial method}

A non-zero polynomial of degree $d$ in one variable has at most $d$ roots. In several variables there is no such bound, since
$x_1x_2$ vanishes on two whole lines, but the right generalisations are just as strong. They turn a combinatorial question into
the question of whether some polynomial can vanish everywhere.

\begin{exercise}[$\star$]\label{pb:schwartz}
\emph{Schwartz--Zippel.} Let $F$ be a field and $S\subset F$ finite.
\begin{enumerate}[label=(\alph*)]
\item Let $f\in F[x_1,\dots,x_n]$ be non-zero of total degree $d$, and $r$ uniform in $S^n$. Show that $P(f(r)=0)\le\frac d{|S|}$,
and that the bound is reached.
\item (Edmonds, 1967.) Let $G$ be a bipartite graph on $\{u_1,\dots,u_n\}\cup\{v_1,\dots,v_n\}$, and $X$ the $n\times n$ matrix with
$X_{ij}=x_{ij}$ (a variable) if $u_iv_j$ is an edge and $0$ otherwise. Show that $\det X\ne0$ in $\Q[x_{ij}]$ if and only if $G$ has
a perfect matching.
\item Deduce a randomised test for perfect matchings that uses one determinant over $\mathbb F_p$ and errs with probability at
most $\frac np$, and only in one direction. Which one?
\end{enumerate}
\end{exercise}

\begin{exercise}[$\star\star$]\label{pb:nullstellensatz}
\emph{Alon's Combinatorial Nullstellensatz.} Let $F$ be a field and $f\in F[x_1,\dots,x_n]$.
\begin{enumerate}[label=(\alph*)]
\item Let $S_1,\dots,S_n\subset F$ be finite. Show that if $\deg_{x_i}f<|S_i|$ for every $i$ and $f$ vanishes on
$S_1\times\dots\times S_n$, then $f=0$.
\item (Alon, 1999.) Suppose that $\deg f=t_1+\dots+t_n$ and that the coefficient of $x_1^{t_1}\cdots x_n^{t_n}$ in $f$ is non-zero.
Show that if $|S_i|>t_i$ for every $i$, then $f(s)\ne0$ for some $s\in S_1\times\dots\times S_n$.
\item (Alon's permanent lemma.) Let $A\in M_n(F)$ have non-zero permanent $\operatorname{per}A=\sum_\sigma\prod_iA_{i\sigma(i)}$,
let $b\in F^n$, and let $S_1,\dots,S_n\subset F$ have two elements each. Show that there is $x\in S_1\times\dots\times S_n$ such
that $(Ax)_i\ne b_i$ for every $i$.
\end{enumerate}
\end{exercise}

\begin{exercise}[$\star\star$]\label{pb:cauchydav}
\emph{Cauchy--Davenport.} Let $p$ be prime and $A,B\subset\Z/p\Z$ non-empty.
\begin{enumerate}[label=(\alph*)]
\item (Cauchy 1813, Davenport 1935.) Prove that $|A+B|\ge\min(p,|A|+|B|-1)$, using Problem~\ref{pb:nullstellensatz}.
\item Deduce $|A_1+\dots+A_k|\ge\min\bigl(p,\sum|A_i|-k+1\bigr)$.
\item Let $a_1,\dots,a_{p-1}$ be non-zero residues, not necessarily distinct. Show that every residue is the sum of some of the
$a_i$. Is $p-1$ the least number that works?
\end{enumerate}
\end{exercise}

\begin{exercise}[$\star\star\star$]\label{pb:erdosheil}
\emph{Sums of distinct elements.} For $A\subset\Z/p\Z$ let $A\dotplus A=\{a+a':a,a'\in A,\ a\ne a'\}$.
\begin{enumerate}[label=(\alph*)]
\item (Dias da Silva and Hamidoune 1994; the proof is due to Alon, Nathanson and Ruzsa, 1995.) Prove the Erdős--Heilbronn
conjecture: $|A\dotplus A|\ge\min(p,2|A|-3)$.
\item Show that the bound is reached for every size $|A|\le\frac{p+3}2$.
\end{enumerate}
\end{exercise}

\begin{exercise}[$\star\star$]\label{pb:egz}
\emph{Erdős--Ginzburg--Ziv (1961).} We show that among any $2n-1$ integers there are $n$ whose sum is divisible by $n$.
\begin{enumerate}[label=(\alph*)]
\item Prove Chevalley--Warning for systems: if $f_1,\dots,f_r\in\mathbb F_p[x_1,\dots,x_m]$ and $\sum_j\deg f_j<m$, then the number
of common zeros of the $f_j$ in $\mathbb F_p^m$ is divisible by $p$.
\item Prove the theorem for $n=p$ prime, with $f_1=\sum_ia_ix_i^{p-1}$ and $f_2=\sum_ix_i^{p-1}$.
\item Show that if the theorem holds for $m$ and for $n$, it holds for $mn$. Conclude.
\item Show that $2n-1$ cannot be replaced by $2n-2$.
\end{enumerate}
\end{exercise}

\begin{exercise}[$\star\star$]\label{pb:hyperplanes}
\emph{Covering a cube minus a point.}
\begin{enumerate}[label=(\alph*)]
\item (Alon and Füredi, 1993.) Show that affine hyperplanes of $\mathbb F_q^n$ that cover every point except $0$ must number at
least $n(q-1)$, and that $n(q-1)$ is enough.
\item (Jamison 1977, Brouwer and Schrijver 1978.) A \emph{blocking set} is a set of points of $\mathbb F_q^n$ that meets every affine
hyperplane. Show that a blocking set has at least $n(q-1)+1$ points. Find one of that size.
\end{enumerate}
\end{exercise}

\begin{exercise}[$\star\star$]\label{pb:kakeya}
\emph{Dvir's Kakeya theorem.} A \emph{Kakeya set} $K\subset\mathbb F_q^n$ contains a line in every direction: for every
$v\ne0$ there is $a$ with $\{a+tv:t\in\mathbb F_q\}\subset K$.
\begin{enumerate}[label=(\alph*)]
\item (Dvir, 2008.) Show that $|K|\ge\binom{q+n-1}n$. In particular $|K|\ge\frac{q^n}{n!}$: a Kakeya set fills a fixed fraction
of the space.
\item For $q$ odd, build a Kakeya set in $\mathbb F_q^2$ with $\frac{q(q+1)}2+\frac{q-1}2$ points.
\end{enumerate}
\end{exercise}

\begin{exercise}[$\star\star$]\label{pb:moore}
\emph{The Moore determinant.} Let $q$ be a prime power and $\Delta(x_1,\dots,x_n)=\det\bigl(x_i^{q^{j-1}}\bigr)_{1\le i,j\le n}$.
\begin{enumerate}[label=(\alph*)]
\item (Moore, 1896.) Show that
\[
\Delta(x_1,\dots,x_n)=\prod_{i=1}^n\ \prod_{c\in\mathbb F_q^{i-1}}\bigl(x_i+c_1x_1+\dots+c_{i-1}x_{i-1}\bigr)
\]
in $\mathbb F_q[x_1,\dots,x_n]$.
\item Deduce that $\alpha_1,\dots,\alpha_n\in\mathbb F_{q^m}$ are linearly independent over $\mathbb F_q$ if and only if
$\Delta(\alpha_1,\dots,\alpha_n)\ne0$.
\item Let $V\subset\mathbb F_{q^m}$ be an $\mathbb F_q$-subspace. Show that $P_V(X)=\prod_{v\in V}(X-v)$ is a linear combination
of $X,X^q,X^{q^2},\dots$, so that $P_V(x+y)=P_V(x)+P_V(y)$.
\end{enumerate}
\end{exercise}

\begin{exercise}[Journey]\label{pb:capset}
\emph{The cap set theorem.} Let $A\subset\mathbb F_3^n$ contain no three distinct points $x,y,z$ with $x+y+z=0$. We prove
$|A|\le3\cdot2.756^n$, following Ellenberg--Gijswijt and Tao.

For a finite set $A$ and a field $F$, a function $T:A^3\to F$ is a \emph{slice} if it has one of the forms $f(x)g(y,z)$,
$f(y)g(x,z)$, $f(z)g(x,y)$. The \emph{slice rank} of $T$ is the least number of slices whose sum is $T$.
\begin{enumerate}[label=(\alph*)]
\item Show that every subspace of $F^A$ of dimension $m$ contains a vector with at least $m$ non-zero coordinates.
\item Let $T(x,y,z)=\sum_{a\in A}c_a\,[x=y=z=a]$ with every $c_a\ne0$. Show that the slice rank of $T$ is $|A|$.
\item Show that for $x,y,z\in A$, $x+y+z=0$ if and only if $x=y=z$. Deduce that $T(x,y,z)=[x+y+z=0]$ has slice rank $|A|$.
\item Write $[x+y+z=0]=\prod_{i=1}^n\bigl(1-(x_i+y_i+z_i)^2\bigr)$ in $\mathbb F_3$, and deduce that the slice rank is at most $3M_n$,
where $M_n$ is the number of $\alpha\in\{0,1,2\}^n$ with $\alpha_1+\dots+\alpha_n\le\frac{2n}3$.
\item Show that $M_n\le\theta^n$ with $\theta=\min_{0<t\le1}\frac{1+t+t^2}{t^{2/3}}$, compute $\theta=2.7551\ldots$, and
conclude.
\end{enumerate}
\end{exercise}

\begin{figure}[ht]
\centering
\fig{ch05-cap}
\caption{A cap of $20$ points in $\mathbb F_3^4$, the largest possible: twenty SET cards with no set among them. The big
$3\times3$ block gives the first two coordinates, the cell inside the block the last two. No three blue cells $x,y,z$
satisfy $x+y+z=0$, even across blocks. The bound of Problem~\ref{pb:capset} is $3M_4=45$ here. It is useless for $n\le3$, but from $n=4$ on it beats the
trivial $3^n$, and by an exponential factor as $n$ grows.}
\label{fig:cap}
\end{figure}

\section{Pseudo-randomness}

The squares modulo $p$, the bits of a linear recurrence, a well-chosen set of $q+1$ residues: these are rigid algebraic objects,
and yet they pass every simple test of randomness. The reason, each time, is a character sum that is as small as a random sum
would be.

\begin{exercise}[$\star\star$]\label{pb:paley}
\emph{Paley graphs.} Let $p\equiv1\pmod4$ be prime. The Paley graph $P_p$ has vertex set $\mathbb F_p$, with $x$ and $y$ joined
when $x-y$ is a non-zero square (this is symmetric because $-1$ is a square).
\begin{enumerate}[label=(\alph*)]
\item Show that two adjacent vertices have $\frac{p-5}4$ common neighbours and two non-adjacent vertices have $\frac{p-1}4$.
Deduce that $P_p$ has $\frac{p(p-1)(p-5)}{48}$ triangles.
\item For $S,T\subset\mathbb F_p$ let $e(S,T)$ be the number of pairs $(s,t)\in S\times T$ with $s,t$ adjacent. Show that
\[
\Bigl|e(S,T)-\frac{|S||T|-|S\cap T|}2\Bigr|\le\frac12\sqrt{p\,|S|\,|T|}.
\]
\item Show that $P_p$ has no clique and no independent set with more than $\sqrt p$ vertices.
\item Deduce from $P_{17}$ that the Ramsey number $R(4,4)$ is at least $18$.
\end{enumerate}
\end{exercise}

\begin{exercise}[$\star\star$]\label{pb:mseq}
\emph{Maximal length sequences.} Let $\alpha$ generate the cyclic group $\mathbb F_{2^n}^*$, $f=x^n+c_{n-1}x^{n-1}+\dots+c_0$ its
minimal polynomial over $\mathbb F_2$, $\operatorname{Tr}(\beta)=\beta+\beta^2+\dots+\beta^{2^{n-1}}$, and
$s_k=\operatorname{Tr}(\alpha^k)\in\mathbb F_2$.
\begin{enumerate}[label=(\alph*)]
\item Show that $s_{k+n}=c_{n-1}s_{k+n-1}+\dots+c_0s_k$ for all $k$.
\item Show that the map $\beta\mapsto\bigl(\operatorname{Tr}(\beta),\operatorname{Tr}(\beta\alpha),\dots,\operatorname{Tr}(\beta\alpha^{n-1})\bigr)$ is a bijection
$\mathbb F_{2^n}\to\mathbb F_2^n$. Deduce that $(s_k)$ has period $2^n-1$, that every non-zero block of $n$ bits appears exactly
once among the windows $(s_k,\dots,s_{k+n-1})$, $0\le k<2^n-1$, and that every non-zero sequence satisfying the recurrence of (a)
is a shift of $(s_k)$.
\item Show that a period contains $2^{n-1}$ ones and $2^{n-1}-1$ zeros, and that for $0<t<2^n-1$,
$\sum_{k=0}^{2^n-2}(-1)^{s_k+s_{k+t}}=-1$.
\end{enumerate}
\end{exercise}

\begin{exercise}[$\star\star\star$]\label{pb:singer}
\emph{Singer's difference sets (1938).} Let $q$ be a prime power, $N=q^2+q+1$, and $g$ a generator of $\mathbb F_{q^3}^*$. View
$\mathbb F_{q^3}$ as a $3$-dimensional vector space over $\mathbb F_q$, fix a $2$-dimensional subspace $U$, and let
$D=\{i\in\Z/N\Z:g^i\in U\}$.
\begin{enumerate}[label=(\alph*)]
\item Show that $D$ is well defined, that $|D|=q+1$, and that every non-zero $d\in\Z/N\Z$ is a difference $i-j$ of elements of
$D$ in exactly one way.
\item Show that the $N$ translates $D+t$ are the lines of a projective plane on the points $\Z/N\Z$: two points lie on exactly one
translate, and two translates meet in exactly one point.
\item Check that $\{0,1,3\}\bmod7$, $\{0,1,3,9\}\bmod13$ and $\{0,1,4,14,16\}\bmod21$ are perfect difference sets.
\end{enumerate}
\end{exercise}

\begin{figure}[ht]
\centering
\fig{ch05-singer}
\caption{The Singer set $\{0,1,3,9\}$ modulo $13$, for $q=3$. Its six chords have the six lengths $\pm1,\dots,\pm6$, so each
non-zero residue is a difference of two of its points exactly once (Problem~\ref{pb:singer}). Turning the picture by $t$ steps
gives the $13$ lines of the projective plane of order $3$.}
\label{fig:singer}
\end{figure}

\begin{exercise}[$\star\star$]\label{pb:sidon}
\emph{Sidon sets.} A set $A$ of integers is a \emph{Sidon set} if all the sums $a+b$, $a\le b$ in $A$, are distinct. Let $F(n)$
be the largest size of a Sidon set in $\{1,\dots,n\}$.
\begin{enumerate}[label=(\alph*)]
\item Show that $F(n)<\sqrt{2n}+1$.
\item (Erdős and Turán, 1941.) Let $p$ be an odd prime. Show that $\{2pk+(k^2\bmod p):0\le k<p\}$ is a Sidon set, with all elements
below $2p^2$.
\item (Erdős and Turán.) Show that $F(n)\le\sqrt n+O(n^{1/4})$. Count the pairs of elements of $A$ that fall in a common window
of $m$ consecutive integers, in two ways.
\item Using Singer sets and the prime number theorem, show that $F(n)\sim\sqrt n$.
\end{enumerate}
\end{exercise}

\begin{exercise}[$\star\star$]\label{pb:randpoly}
\emph{Random polynomials, random permutations.} Let $f$ be a uniformly random monic polynomial of degree $n$ over $\mathbb F_q$,
and $M(j,q)$ the number of monic irreducible polynomials of degree $j$.
\begin{enumerate}[label=(\alph*)]
\item Show that, for $n\ge2$, $f$ is squarefree with probability exactly $1-\frac1q$.
\item Show that $f$ is irreducible with probability $\frac1n+\varepsilon$ with $|\varepsilon|\le\frac2n q^{-n/2}$.
\item For a partition $\lambda$ of $n$ with $m_j$ parts equal to $j$, show that the probability that $f$ factors into irreducibles of
degrees given by $\lambda$ (with multiplicity) tends, as $q\to\infty$, to $1/z_\lambda$, where $z_\lambda=\prod_jj^{m_j}m_j!$.
Interpret $1/z_\lambda$.
\end{enumerate}
\end{exercise}

\begin{exercise}[$\star\star$]\label{pb:noroots}
\emph{Roots of a random polynomial.} Let $f$ be a uniformly random monic polynomial of degree $n\ge1$ over $\mathbb F_q$ and $R$ its
number of roots in $\mathbb F_q$.
\begin{enumerate}[label=(\alph*)]
\item Show that if $n\ge q$, the values $f(a)$, $a\in\mathbb F_q$, are independent and uniform, so that $R$ is exactly binomial
with parameters $q$ and $\frac1q$. In particular $P(R=0)=(1-\frac1q)^q$.
\item For $n<q$, show that
\[
P(R=0)=\sum_{k=0}^n(-1)^k\binom qk q^{-k},
\]
and compare with the probability $\sum_{k=0}^n\frac{(-1)^k}{k!}$ that a random permutation of $n$ points has no fixed point.
\end{enumerate}
\end{exercise}

\begin{exercise}[$\star\star\star$]\label{pb:waring}
\emph{Two $k$-th powers.} Let $p$ be prime, $k\ge2$ a divisor of $p-1$, and $N(c)$ the number of $(x,y)\in\mathbb F_p^2$ with
$x^k+y^k=c$.
\begin{enumerate}[label=(\alph*)]
\item Show that for $c\ne0$,
\[
|N(c)-p|\le(k-1)(k-2)\sqrt p+k-1.
\]
\item Deduce that every residue modulo $p$ is a sum of two $k$-th powers when $p\ge k^4$.
\item What happens when $k$ does not divide $p-1$?
\end{enumerate}
\end{exercise}

\section{Hints}

\begin{hint}{pb:conics}
(a) Complete the square: $4a(ax^2+bx+c)=(2ax+b)^2-D$. Then count the solutions of $y^2-z^2=D$. (b) For fixed $y$, the number of
$x$ is $1+\chi(\text{discriminant})$.
\end{hint}

\begin{hint}{pb:sumsquares}
Detect the equation with $\frac1p\sum_te(t(\sum x_i^2-c))$, and evaluate $\sum_xe(tx^2)$ with Tome~1, Exercise~6.7.
\end{hint}

\begin{hint}{pb:fermatcubic}
(a) The number of cube roots of $u$ is $1+\chi(u)+\chi^2(u)$. (c) Raise $g(\chi)$ to the cube modulo $3$: in characteristic $3$,
cubing is additive.
\end{hint}

\begin{hint}{pb:nilpotent}
(b) A matrix is the same thing as a decomposition $V=U\oplus W$, a nilpotent map on $U$ and an invertible map on $W$. (c) Show
by induction that $\sum_{k\le n}q^{\binom k2}/D_k=q^{\binom{n+1}2}/D_n$, where $D_k=\prod_{i=1}^k(q^i-1)$.
\end{hint}

\begin{hint}{pb:rankdist}
(a) Choose the image first. (b) Pull $q^{n}$ out of every factor. (c) The probabilities add up to $1$; justify the limit with a
uniform bound.
\end{hint}

\begin{hint}{pb:flags}
(b) A complement is the graph of a linear map from a fixed complement to $U$.
\end{hint}

\begin{hint}{pb:schwartz}
(a) Induction on $n$. Write $f$ as a polynomial in $x_1$ whose coefficients are polynomials in the other variables:
\[
f=\sum_{i\le k}x_1^i\,g_i(x_2,\dots,x_n),\qquad g_k\ne0.
\]
\end{hint}

\begin{hint}{pb:nullstellensatz}
(b) Shrink the $S_i$ to size $t_i+1$, and divide $f$ by the polynomials $\prod_{s\in S_i}(x_i-s)$. If $f$ vanished on the grid, the
remainder would vanish by (a).
\end{hint}

\begin{hint}{pb:cauchydav}
(a) If $A+B$ were inside a set $C$ of size $|A|+|B|-2$, the polynomial $\prod_{c\in C}(x+y-c)$ would vanish on $A\times B$.
\end{hint}

\begin{hint}{pb:erdosheil}
Use $(x-y)\prod_{c\in C}(x+y-c)$ and the monomial $x^{k-1}y^{k-2}$. Its coefficient is a Catalan number.
\end{hint}

\begin{hint}{pb:egz}
(a) The number of common zeros is $\sum_x\prod_j(1-f_j(x)^{p-1})$ modulo $p$. (c) Extract groups of $n$ one at a time.
\end{hint}

\begin{hint}{pb:hyperplanes}
(a) Compare $\prod_i(1-a_i\cdot x)$ with $\prod_j\prod_{s\ne0}(x_j-s)$. (b) The hyperplanes $\{x:a\cdot x=1\}$ avoid $0$; swap the
roles of $a$ and $x$.
\end{hint}

\begin{hint}{pb:kakeya}
(a) A non-zero polynomial of degree at most $q-1$ vanishing on $K$ exists if $K$ is small. Restrict it to a line and look at the
top coefficient. (b) Use the lines $y=mx+m^2$.
\end{hint}

\begin{hint}{pb:moore}
(a) Expand along the last row, and note that the last row is a combination of the others when $x_n$ is an $\mathbb F_q$-combination
of $x_1,\dots,x_{n-1}$.
\end{hint}

\begin{hint}{pb:capset}
(b) Contract $T$ against a function $h(z)$ that is orthogonal to the $f$'s of the third kind of slices and has large support.
(d) In each monomial one of the three variables has degree at most $\frac{2n}3$.
\end{hint}

\begin{hint}{pb:paley}
(a) Count with $\frac{1+\chi}2$ and use Problem~\ref{pb:conics}(a). (b) Cauchy--Schwarz in $s$, then sum over all $s\in\mathbb F_p$.
(c) If $K$ is a clique and $n$ a non-square, $(c,c')\mapsto c+nc'$ is injective on $K\times K$.
\end{hint}

\begin{hint}{pb:mseq}
(b) The trace is a non-zero linear form, and $(\beta,\gamma)\mapsto\operatorname{Tr}(\beta\gamma)$ is non-degenerate.
(c) $s_k+s_{k+t}=\operatorname{Tr}(\alpha^k(1+\alpha^t))$.
\end{hint}

\begin{hint}{pb:singer}
(a) $i-j=d$ with $i,j\in D$ means that the point $g^j$ lies in $U\cap g^{-d}U$. Why are these two planes different?
\end{hint}

\begin{hint}{pb:sidon}
(b) Equal sums give equal $k+l$ and equal $k^2+l^2$ modulo $p$. (c) Window length $m\approx n^{3/4}$; a difference $d<m$ occurs at
most once.
\end{hint}

\begin{hint}{pb:randpoly}
(a) $f=g^2h$ with $h$ squarefree, in generating functions. (c) Choose $m_j$ irreducibles of degree $j$ with repetition.
\end{hint}

\begin{hint}{pb:noroots}
(a) Evaluation at all points of $\mathbb F_q$ is a surjective linear map on polynomials of degree $<q$. (b) Inclusion--exclusion.
\end{hint}

\begin{hint}{pb:waring}
Write the number of $k$-th roots of $u$ as a sum of characters, and use Jacobi sums of absolute value $\sqrt p$.
\end{hint}

\section{Solutions}

\begin{solution}{pb:conics}
\idea A square completion reduces every quadratic character sum to the one for $y^2-D$, and $y^2-z^2=D$ factors as
$(y-z)(y+z)=D$, which is easy to count.

\writeup (a) Since $p$ is odd, $4a\ne0$ and $\chi(4a)=\chi(a)$. From $4a(ax^2+bx+c)=(2ax+b)^2-D$ we get
$\chi(ax^2+bx+c)=\chi(a)\chi\bigl((2ax+b)^2-D\bigr)$, and $y=2ax+b$ runs over $\mathbb F_p$ once when $x$ does. So the sum is
$\chi(a)\sum_y\chi(y^2-D)$. If $D=0$ this is $\chi(a)(p-1)$. If $D\ne0$, the number of $z$ with $z^2=w$ is $1+\chi(w)$, so
\[
\sum_y\chi(y^2-D)=\#\{(y,z):y^2-z^2=D\}-p.
\]
With $u=y-z$, $v=y+z$ (an invertible change of variables, $2$ being invertible) the equation is $uv=D$, which has $p-1$
solutions. So the sum is $-1$, and the total is $-\chi(a)$.

(b) Suppose first $a\ne0$. For fixed $y$, the equation $ax^2+(by)x+(cy^2-1)=0$ in $x$ has $1+\chi(\Delta_y)$ solutions, where
$\Delta_y=b^2y^2-4a(cy^2-1)=Dy^2+4a$. The polynomial $Dy^2+4a$ has leading coefficient $D\ne0$ and discriminant $-16aD\ne0$, so
by (a)
\[
N=\sum_y\bigl(1+\chi(Dy^2+4a)\bigr)=p-\chi(D).
\]
If $a=0$, then $D=b^2\ne0$ and $\chi(D)=1$. The equation is $y(bx+cy)=1$: any $y\ne0$, and then $x=(y^{-1}-cy)/b$ is unique. That is
$p-1=p-\chi(D)$ solutions.

(c) For $n\ne0$, $Q=n$ is the equation $Q'=1$ with $Q'=Q/n$, whose discriminant $D/n^2$ has the same character. So $Q=n$ has
$p-\chi(D)\ge p-1>0$ solutions. For $x^2+y^2$, $D=-4$ and $\chi(-4)=\chi(-1)$.

\tools Tome~1, Chapter~6 (Legendre symbol).
\leads $x^2+y^2=1$ is a circle with $p+1$ or $p-1$ points: the $p-1$ case corresponds to a circle with two points at
infinity, $x\pm iy=0$, defined over $\mathbb F_p$. Every smooth conic has $p+1$ points in the projective plane, and the problem
is the affine shadow of that fact.
\end{solution}

\begin{solution}{pb:sumsquares}
\idea Additive characters detect the equation, and the sum over each coordinate is a quadratic Gauss sum, whose square is
$\pm p$. All the answer needs is the sign.

\writeup Let $G=\sum_t\chi(t)e(t)$. By Tome~1, Exercise~6.7, $G^2=\chi(-1)p$. For $t\ne0$,
\[
S(t)=\sum_xe(tx^2)=\sum_u(1+\chi(u))e(tu)=\sum_u\chi(u)e(tu)=\chi(t)G,
\]
using $\sum_ue(tu)=0$ and the substitution $u\mapsto u/t$. Since $\frac1p\sum_te(tm)$ is $1$ if $m=0$ and $0$ otherwise,
\[
N_n(c)=\frac1p\sum_{t}e(-tc)S(t)^n=p^{n-1}+\frac{G^n}p\sum_{t\ne0}\chi(t)^ne(-tc).
\]
(a) If $n$ is even, $\chi(t)^n=1$, $G^n=(\chi(-1)p)^{n/2}=\eta\bigl((-1)^{n/2}\bigr)p^{n/2}$, and $\sum_{t\ne0}e(-tc)$ is $p-1$ if
$c=0$ and $-1$ otherwise. So
\[
N_n(c)=p^{n-1}+p^{n/2-1}\eta\bigl((-1)^{n/2}\bigr)\cdot\begin{cases}p-1&c=0,\\-1&c\ne0.\end{cases}
\]
If $n$ is odd, $\chi(t)^n=\chi(t)$. For $c=0$ the sum $\sum_{t\ne0}\chi(t)$ vanishes and $N_n(0)=p^{n-1}$. For $c\ne0$,
$\sum_t\chi(t)e(-tc)=\chi(-c)G$ (substitute $t\mapsto-t/c$), so
\[
N_n(c)=p^{n-1}+\frac{\chi(-c)G^{n+1}}p=p^{n-1}+\chi(-c)\chi(-1)^{\frac{n+1}2}p^{\frac{n-1}2}=p^{n-1}+\eta\bigl((-1)^{\frac{n-1}2}c\bigr)p^{\frac{n-1}2},
\]
since $\chi(-1)^{\frac{n+1}2+1}=\chi(-1)^{\frac{n-1}2}$.

(b) Let $P(c)=N_n(c)/p^n$. The distance is $\frac12\sum_c|P(c)-\frac1p|$. For $n$ even, $|P(0)-\frac1p|=(p-1)p^{-n/2-1}$ and
$|P(c)-\frac1p|=p^{-n/2-1}$ for the $p-1$ others, so the distance is $(p-1)p^{-n/2-1}=(1-\frac1p)p^{-n/2}$. For $n$ odd,
$P(0)=\frac1p$ and $|P(c)-\frac1p|=p^{-(n+1)/2}$ for $c\ne0$, giving $\frac{p-1}2p^{-(n+1)/2}$.

\tools Tome~1, Chapter~6 (Exercise~6.7, the quadratic Gauss sum); Tome~3, Chapter~8 (total variation).
\leads So the sum of $n$ random squares is uniform up to an error $p^{-n/2}$, exactly the size that the square root
cancellation of Gauss sums predicts. For $n=4$ and $c=m\bmod p$ the formula $p^3-p$ is the local factor in Jacobi's count of
representations of $m$ as a sum of four squares, which Part~II will meet with modular forms.
\end{solution}

\begin{solution}{pb:fermatcubic}
\idea Counting solutions of $u+v=1$ with $u,v$ cubes is a sum of Jacobi sums. The only non-trivial one is $J(\chi,\chi)$, an
integer of $\Z[\omega]$ of norm $p$, and a congruence modulo $3$ picks it out among the six associates.

\writeup (a) Since $p\equiv1\pmod 3$, $\mathbb F_p^*$ is cyclic of order divisible by $3$; it has two characters $\chi,\chi^2=\bar\chi$
of order $3$. For $u\ne0$, the number of $x$ with $x^3=u$ is $3$ if $u$ is a cube and $0$ otherwise, which is $1+\chi(u)+\chi^2(u)$;
for $u=0$ it is $1$, which is also $1+\chi(0)+\chi^2(0)$ with our conventions. Writing $\varepsilon$ for the trivial character,
\[
N=\sum_{u+v=1}\Bigl(\sum_{i=0}^2\chi^i(u)\Bigr)\Bigl(\sum_{j=0}^2\chi^j(v)\Bigr)=\sum_{i,j}J(\chi^i,\chi^j).
\]
Now $J(\varepsilon,\varepsilon)=p$ and $J(\varepsilon,\psi)=\sum_v\psi(v)=0$ for $\psi$ non-trivial. Next
$J(\chi,\bar\chi)=\sum_{u\ne0,1}\chi\bigl(\frac u{1-u}\bigr)$, and $u\mapsto\frac u{1-u}$ is a bijection from
$\mathbb F_p\setminus\{0,1\}$ to $\mathbb F_p\setminus\{0,-1\}$; so $J(\chi,\bar\chi)=-\chi(-1)=-1$, because
$\chi(-1)^2=\chi(1)=1$ and $\chi(-1)^3=1$ force $\chi(-1)=1$. The same holds for $J(\bar\chi,\chi)$. Finally
$J(\bar\chi,\bar\chi)=\overline{J(\chi,\chi)}$. So $N=p-2+2\operatorname{Re}J(\chi,\chi)$.

(b) Expanding,
\[
g(\alpha)g(\beta)=\sum_s e(s)\sum_{x+y=s}\alpha(x)\beta(y).
\]
For $s\ne0$, put $x=su$, $y=sv$: the inner sum is $\alpha\beta(s)J(\alpha,\beta)$. For $s=0$ it is
\[
\sum_x\alpha(x)\beta(-x)=\beta(-1)\sum_x\alpha\beta(x)=0,
\]
since $\alpha\beta$ is non-trivial. So
\[
g(\alpha)g(\beta)=J(\alpha,\beta)\sum_s\alpha\beta(s)e(s)=J(\alpha,\beta)g(\alpha\beta).
\]

Next, $|g(\chi)|^2=\sum_{x,y\ne0}\chi(x)\bar\chi(y)e(x-y)$. Put $x=ty$: this is
$\sum_{t\ne0}\chi(t)\sum_{y\ne0}e(y(t-1))$. The inner sum is $p-1$ for $t=1$ and $-1$ otherwise, so
$|g|^2=(p-1)-\sum_{t\ne0,1}\chi(t)=(p-1)+1=p$. Also $\overline{g(\chi)}=\sum\bar\chi(t)e(-t)=\chi(-1)g(\bar\chi)=g(\bar\chi)$.
Apply the first identity twice: $g(\chi)^2=J(\chi,\chi)g(\chi^2)=J(\chi,\chi)g(\bar\chi)$, so
$g(\chi)^3=J(\chi,\chi)g(\bar\chi)g(\chi)=J(\chi,\chi)|g(\chi)|^2=pJ(\chi,\chi)$.

(c) Work in the ring $R=\Z[\omega,\zeta]$ with $\zeta=e^{2i\pi/p}$. Modulo $3R$, $(a+b)^3\equiv a^3+b^3$ (the middle binomial
coefficients are $3$), so
\[
g(\chi)^3\equiv\sum_t\chi(t)^3\zeta^{3t}=\sum_{t\ne0}\zeta^{3t}=-1\pmod{3R},
\]
because $\chi(t)^3=1$ for $t\ne0$ and $3t$ runs over $\mathbb F_p^*$. Since $p\equiv1\pmod3$, (b) gives $J(\chi,\chi)\equiv-1$,
that is $\frac{J+1}3\in R$. But $\frac{J+1}3=\frac{a'+1}3+\frac{b'}3\omega$ is in $\Q(\omega)$ and is an algebraic integer, so it lies
in $\Z[\omega]$; as $1,\omega$ is a basis of $\Z[\omega]$ and of $\Q(\omega)$, $3\mid a'+1$ and $3\mid b'$.

(d) Let $a=2\operatorname{Re}J=2a'-b'$ (because $\operatorname{Re}\omega=-\frac12$). Then $a\equiv-2\equiv1\pmod 3$. By (b),
$|J|^2=\frac{|g|^6}{p^2}=p$, and $|a'+b'\omega|^2=a'^2-a'b'+b'^2$, so $4p=(2a'-b')^2+3b'^2=a^2+27b^2$ with $b=\frac{b'}3$. And
$N=p-2+a$. Since $27b^2\ge0$, $|a|\le2\sqrt p$.

For $p=7$: $28=1+27$, $a=1$, $N=6$. Directly: the cubes mod $7$ are $0,1,6$, and $u+v=1$ with $u,v$ cubes forces $\{u,v\}=\{0,1\}$,
giving $1\cdot3+3\cdot1=6$ solutions. For $p=13$: $52=25+27$, $a\equiv1\pmod3$ picks $a=-5$, so $N=6$. (Why is $a$ determined?
If $4p=a^2+27b^2=c^2+27d^2$, the numbers $\frac{a\pm3b\sqrt{-3}}2$ and $\frac{c\pm3d\sqrt{-3}}2$ have norm $p$ in the unique
factorisation ring $\Z[\omega]$, so they are associates, and the congruence modulo $3$ leaves only $c=a$.)

\tools Tome~1, Chapters~5 and~6 (characters, Gauss sums); Tome~1, Chapters~10 and~11 ($\Z[\omega]$, algebraic integers); Tome~2, Chapter~9 (Frobenius).
\leads This is Gauss's result in article~358 of the \emph{Disquisitiones} (1801), and André Weil's 1949 paper on equations over
finite fields starts from it: the projective curve $x^3+y^3=z^3$ has $p+1-a$ points, and $|a|\le2\sqrt p$ is the Hasse bound for
this elliptic curve. Figure~\ref{fig:cubic} shows $\frac a{2\sqrt p}$; because the curve has complex multiplication by $\Z[\omega]$,
its distribution is not the Sato--Tate semicircle but the arcsine-like law of $\cos\theta$ for uniform $\theta$, which piles up at
$\pm1$.
\end{solution}

\begin{solution}{pb:nilpotent}
\idea Every matrix splits uniquely into a nilpotent part and an invertible part on complementary subspaces. Counting matrices
this way gives one identity per $n$, and these identities determine all the $N_k$, so it is enough to check that $q^{k^2-k}$
satisfies them.

\writeup (a) The kernels $\ker A^k$ increase and the images $\operatorname{im}A^k$ decrease, and both stabilise by $k=n$ (the
dimension cannot change more than $n$ times). If $x\in\ker A^n\cap\operatorname{im}A^n$, write $x=A^ny$; then $A^{2n}y=0$, so
$y\in\ker A^{2n}=\ker A^n$ and $x=0$. By rank--nullity the dimensions add up to $n$, so $V=\ker A^n\oplus\operatorname{im}A^n$.
Both are $A$-stable. $A$ is nilpotent on $\ker A^n$, and on $\operatorname{im}A^n$ it is onto
($A(\operatorname{im}A^n)=\operatorname{im}A^{n+1}=\operatorname{im}A^n$), hence invertible. If $V=U\oplus W$ is another such
decomposition, then $A^n=0$ on $U$, so $U\subset\ker A^n$, and $W=A^nW\subset\operatorname{im}A^n$; the dimensions force equality.

(b) By (a), $A\mapsto(U,W,A|_U,A|_W)$ is a bijection from $M_n(\mathbb F_q)$ to the set of quadruples where $V=U\oplus W$, $A|_U$
is nilpotent on $U$ and $A|_W$ is invertible on $W$. $GL_n$ acts transitively on the ordered pairs $(U,W)$ with
$\dim U=k$, $V=U\oplus W$ (send a basis adapted to one pair to a basis adapted to another), with stabiliser $GL(U)\times GL(W)$.
So there are $\frac{|GL_n|}{|GL_k||GL_{n-k}|}$ such pairs, and
\[
q^{n^2}=\sum_{k=0}^n\frac{|GL_n|}{|GL_k|\,|GL_{n-k}|}\,N_k\,|GL_{n-k}|=|GL_n|\sum_{k=0}^n\frac{N_k}{|GL_k|}.
\]

(c) The identity of (b) for $n$ determines $N_n$ from $N_0,\dots,N_{n-1}$, so it suffices to show that $N_k=q^{k^2-k}$ satisfies
it. Write $|GL_k|=\prod_{i=0}^{k-1}(q^k-q^i)=q^{\binom k2}D_k$ with $D_k=\prod_{i=1}^k(q^i-1)$. Then
$\frac{q^{k^2-k}}{|GL_k|}=\frac{q^{\binom k2}}{D_k}$ and $\frac{q^{n^2}}{|GL_n|}=\frac{q^{\binom{n+1}2}}{D_n}$, so we need
\[
\sum_{k=0}^n\frac{q^{\binom k2}}{D_k}=\frac{q^{\binom{n+1}2}}{D_n}.
\]
For $n=0$ both sides are $1$. If it holds for $n$, adding the term $k=n+1$ gives
$\frac{q^{\binom{n+1}2}}{D_n}+\frac{q^{\binom{n+1}2}}{D_{n+1}}=\frac{q^{\binom{n+1}2}(q^{n+1}-1+1)}{D_{n+1}}=\frac{q^{\binom{n+2}2}}{D_{n+1}}$.

(d) $A\mapsto I+A$ is a bijection from nilpotent to unipotent matrices, so there are $q^{n^2-n}$ unipotent matrices.

\tools Tome~0, Chapter~7 (rank--nullity, invariant subspaces); Tome~2, Chapter~3 (orbit--stabiliser); Companion~C1, Chapter~3 for the
order of $GL_n(\mathbb F_q)$.
\leads $q^{n^2-n}=q^{2\binom n2}$ is the square of the order of a Sylow $p$-subgroup of $GL_n$ (the unitriangular matrices).
Steinberg proved that in every finite group of Lie type the number of unipotent elements is the square of the order of a Sylow
$p$-subgroup. Over a general field the nilpotent matrices form the \emph{nilpotent cone}, of dimension $n^2-n$.
\end{solution}

\begin{solution}{pb:rankdist}
\idea A matrix of rank $r$ is a choice of an $r$-dimensional image and of a surjection onto it. In the limit only the corank
matters, and summing the limit law to $1$ is a partition identity.

\writeup (a) A matrix $A$ of rank $r$ is a linear map $\mathbb F_q^n\to\mathbb F_q^n$ with image $W$ of dimension $r$. Given $W$, with a
basis, the maps onto $W$ correspond to $r\times n$ matrices of rank $r$, whose rows are independent: $\prod_{i<r}(q^n-q^i)$ of them.
The number of $W$ is the Gaussian binomial $\prod_{i<r}\frac{q^n-q^i}{q^r-q^i}$ (Companion~C1, Theorem~3.18). Multiply.

(b) Let $r=n-k$. Factor $q^n$ out of each $q^n-q^i$ and $q^r$ out of each $q^r-q^i$:
\[
\frac{\#\{\text{rank }r\}}{q^{n^2}}=q^{2nr-r^2-n^2}\,\frac{\prod_{i<r}(1-q^{i-n})^2}{\prod_{i<r}(1-q^{i-r})}
=q^{-k^2}\,\frac{\prod_{j=k+1}^{n}(1-q^{-j})^2}{\prod_{j=1}^{n-k}(1-q^{-j})}.
\]
As $n\to\infty$ the products converge (since $\sum q^{-j}<\infty$), to $\prod_{j>k}(1-q^{-j})^2/\prod_{j\ge1}(1-q^{-j})$, which is the
stated limit.

(c) Let $c_q=\prod_{j\ge1}(1-q^{-j})>0$. The $n$-th probabilities are at most $q^{-k^2}/c_q$, a summable bound, and they add up to
$1$ for each $n$; by dominated convergence the limits add up to $1$:
\[
\sum_{k\ge0}q^{-k^2}\frac{\prod_{j>k}(1-q^{-j})}{\prod_{j\le k}(1-q^{-j})}=1,\quad\text{that is}\quad
\sum_k\frac{x^{k^2}}{(x)_k^2}=\frac1{\prod_j(1-x^j)}\ \text{ at }x=\tfrac1q,
\]
after dividing by $c_q$ and writing $\prod_{j>k}(1-x^j)=c_q/(x)_k$. Both sides are power series in $x$ with non-negative integer
coefficients, convergent for $|x|<1$, and they agree at $x=\frac1q$ for every prime power $q$, a sequence tending to $0$. By the
identity theorem for power series (Tome~4, Chapter~9) they are equal.

\tools Companion~C1, Chapter~3 (Gaussian binomials; Exercise~3.11 is the case $k=0$); Tome~3, Chapter~5 (dominated convergence);
Tome~4, Chapter~9 (power series).
\leads A partition has a largest square $k\times k$ in its Young diagram (its Durfee square), and the two pieces outside the square
are partitions into parts at most $k$: this is the combinatorial proof, $\frac{x^{k^2}}{(x)_k^2}$ counting partitions with
Durfee square $k$. For $q=2$ the limit law is $0.2888$, $0.5776$, $0.1284$, $0.0052$ for coranks $0,1,2,3$, the Cohen--Lenstra
type law that also governs, conjecturally, the class groups of number fields.
\end{solution}

\begin{solution}{pb:flags}
\idea Build a flag one step at a time, and count complements as graphs of linear maps.

\writeup (a) $V$ has $\frac{q^n-1}{q-1}=[n]_q$ lines. A complete flag is a line $V_1$ followed by a complete flag of $V/V_1$, a space
of dimension $n-1$ (subspaces of $V/V_1$ are subspaces of $V$ containing $V_1$). By induction there are $[n]_q[n-1]_q\cdots[1]_q$.

(b) Fix one complement $W_0$ and let $\pi:V\to W_0$ be the projection along $U$. A complement $W$ meets $U$ trivially and has
dimension $n-k$, so $\pi|_W$ is an isomorphism $W\to W_0$, and $W=\{w+f(w):w\in W_0\}$ where $f=(\mathrm{id}-\pi)\circ(\pi|_W)^{-1}$ maps
$W_0$ to $U$. Conversely, every linear $f:W_0\to U$ gives such a complement, the graph of $f$. So the complements correspond to
$\operatorname{Hom}(W_0,U)$, which has $q^{k(n-k)}$ elements.

(c) $GL_n$ acts transitively on complete flags (choose a basis $e_1,\dots,e_n$ with $V_i=\langle e_1,\dots,e_i\rangle$), and the
stabiliser of the standard flag is the group $B$ of invertible upper triangular matrices: $q-1$ choices for each diagonal entry
and $q$ for each of the $\binom n2$ entries above it. So $|GL_n|=[n]_q!\,(q-1)^nq^{\binom n2}$.

(d) At $q=1$, $[n]_q!$ becomes $n!$, the number of orderings of an $n$-element set, which is what a ``complete flag of a set''
$\emptyset\subset S_1\subset\dots\subset S_n$ is. And $q^{k(n-k)}$ becomes $1$: a subset has exactly one complement. In (c),
$|GL_n|/(q-1)^n=q^{\binom n2}[n]_q!$ becomes $n!=|\mathfrak S_n|$: the symmetric group behaves like $GL_n$ over a ``field with
one element''.

\tools Tome~0, Chapter~7; Tome~2, Chapter~3; Companion~C1, Chapter~3 ($q$-analogues).
\leads Tits suggested in 1957 that $\mathfrak S_n$ should be $GL_n(\mathbb F_1)$, and the dream of a geometry over $\mathbb F_1$ that
would explain the Riemann hypothesis by imitating Weil's proof over $\mathbb F_q$ is still alive. Summing over the cells of the
Bruhat decomposition also gives $[n]_q!=\sum_{\sigma\in\mathfrak S_n}q^{\operatorname{inv}\sigma}$.
\end{solution}

\begin{solution}{pb:schwartz}
\idea A polynomial in $x_1$ with coefficients in the other variables can only vanish if either its top coefficient vanishes
(probability by induction) or $x_1$ hits one of its few roots.

\writeup (a) Induction on $n$. For $n=1$, $f$ has at most $d$ roots. For $n\ge2$, write $f=\sum_{i=0}^kx_1^ig_i(x_2,\dots,x_n)$ with
$g_k\ne0$; then $\deg g_k\le d-k$. Let $r=(r_1,r')$. By induction $P(g_k(r')=0)\le\frac{d-k}{|S|}$. When $g_k(r')\ne0$, the polynomial
$f(x_1,r')$ is non-zero of degree $k$, so for each such $r'$, at most $k$ values of $r_1$ (independent of $r'$) make it vanish.
So $P(f(r)=0)\le\frac{d-k}{|S|}+\frac k{|S|}$. For $f=\prod_{i=1}^d(x_1-s_i)$ with distinct $s_i\in S$ the probability is exactly
$\frac d{|S|}$.

(b) $\det X=\sum_\sigma\operatorname{sgn}\sigma\prod_iX_{i\sigma(i)}$. The term of $\sigma$ is non-zero exactly when every
$u_iv_{\sigma(i)}$ is an edge, that is when $\sigma$ is a perfect matching, and it is then the monomial $\pm\prod_ix_{i\sigma(i)}$.
Different $\sigma$ give different monomials, so nothing cancels: $\det X\ne0$ if and only if a perfect matching exists.

(c) Substitute independent uniform values of $\mathbb F_p$ for the $x_{ij}$ and compute the determinant modulo $p$. If there is no
perfect matching, $\det X$ is the zero polynomial and the answer is always $0$. If there is one, $\det X$ reduced mod $p$ is still
non-zero (its coefficients are $\pm1$) of degree $n$, and by (a) it vanishes with probability at most $\frac np$. So the test
says ``matching'' only when there is one: it can only miss matchings.

\tools Tome~0, Chapter~3 (roots of polynomials); Tome~0, Chapter~7 (determinants).
\leads Lovász (1979) extended the test to general graphs with the Tutte matrix, and Mulmuley, Vazirani and Vazirani (1987) used
the same determinant, with an isolation lemma, to find matchings in parallel. Polynomial identity testing is the most famous
problem that is easy with randomness and not known to be easy without it.
\end{solution}

\begin{solution}{pb:nullstellensatz}
\idea Dividing $f$ by the polynomials $g_i=\prod_{s\in S_i}(x_i-s)$ leaves a remainder of low degree in each variable that agrees
with $f$ on the grid. If $f$ vanished on the grid, the remainder would be $0$, and every top-degree monomial of $f$ would be
divisible by some $x_i^{t_i+1}$.

\writeup (a) Induction on $n$. For $n=1$ a polynomial of degree $<|S_1|$ with $|S_1|$ roots is $0$. For $n\ge2$ write
$f=\sum_{i<|S_n|}f_i(x_1,\dots,x_{n-1})x_n^i$. For fixed $s'\in S_1\times\dots\times S_{n-1}$, the polynomial $f(s',x_n)$ has degree
$<|S_n|$ and $|S_n|$ roots, so all $f_i(s')=0$. By induction each $f_i$ is $0$.

(b) Replacing $S_i$ by a subset, we may assume $|S_i|=t_i+1$. Suppose $f$ vanishes on the grid. Let
$g_i=\prod_{s\in S_i}(x_i-s)=x_i^{t_i+1}-h_i(x_i)$ with $\deg h_i\le t_i$. As long as some monomial of $f$ contains $x_i^{t_i+1}$,
replace that factor by $h_i(x_i)$. Each replacement subtracts $g_i\cdot m$ for a monomial $m$, does not raise the total degree, and
lowers the degree in $x_i$; the process stops at $\bar f$ with $\deg_{x_i}\bar f\le t_i$ and
\[
f=\bar f+\sum_ig_iu_i,\qquad\deg(g_iu_i)\le\deg f.
\]
Since $g_i=0$ on the grid, $\bar f=f=0$ there, so $\bar f=0$ by (a). Now look at the part of degree $d=\deg f$: on the right it is
$\sum_ix_i^{t_i+1}\cdot(\text{the part of degree }d-t_i-1\text{ of }u_i)$, a sum of monomials each divisible by some
$x_i^{t_i+1}$. But $x_1^{t_1}\cdots x_n^{t_n}$ has degree $d$, is divisible by no $x_i^{t_i+1}$, and has a non-zero coefficient in $f$.
Contradiction.

(c) Let $f(x)=\prod_{i=1}^n\bigl((Ax)_i-b_i\bigr)$, of degree $n$. Its part of degree $n$ is $\prod_i\sum_jA_{ij}x_j$, in which the
coefficient of $x_1x_2\cdots x_n$ is $\sum_\sigma\prod_iA_{i\sigma(i)}=\operatorname{per}A\ne0$. Apply (b) with all $t_i=1$.

\tools Tome~0, Chapter~3 (division, roots).
\leads Alon's paper (1999) derives a dozen classical theorems from (b), several of which follow here. The permanent lemma gives,
for instance, that for every $n\times n$ matrix with non-zero permanent over $\mathbb F_p$ and every $b$ there is a $0$--$1$ vector
$x$ with $Ax$ differing from $b$ in every coordinate (Alon and Tarsi's conjecture on nowhere-zero points for $p$ odd builds on it).
\end{solution}

\begin{solution}{pb:cauchydav}
\idea If the sumset were small, a polynomial of controlled degree would vanish on $A\times B$; its key coefficient is a binomial
coefficient below $p$, hence non-zero.

\writeup (a) If $|A|+|B|>p$, then for every $c$ the sets $A$ and $c-B$ have total size $>p$, so they meet: $c\in A+B$. Suppose
$|A|+|B|\le p$ and $|A+B|\le|A|+|B|-2$, and choose $C\supset A+B$ with $|C|=|A|+|B|-2$. The polynomial
$f(x,y)=\prod_{c\in C}(x+y-c)$ vanishes on $A\times B$ and has degree $|A|+|B|-2$. Its top part is $(x+y)^{|A|+|B|-2}$, in which
the coefficient of $x^{|A|-1}y^{|B|-1}$ is $\binom{|A|+|B|-2}{|A|-1}$, not divisible by $p$ since $|A|+|B|-2<p$. Problem~\ref{pb:nullstellensatz}(b)
with $t_1=|A|-1<|A|$ and $t_2=|B|-1<|B|$ gives $(a,b)\in A\times B$ with $f(a,b)\ne0$, that is $a+b\notin C$: a contradiction.

(b) Induction on $k$: $|A_1+\dots+A_k|\ge\min\bigl(p,|A_1+\dots+A_{k-1}|+|A_k|-1\bigr)$, and if the first set already has $p$
elements so does the sum.

(c) Let $A_i=\{0,a_i\}$, of size $2$. By (b), the set of subset sums $A_1+\dots+A_{p-1}$ has at least
$\min(p,2(p-1)-(p-1)+1)=p$ elements. With $p-2$ copies of $1$, the residue $p-1$ is not a subset sum, so $p-1$ terms are needed.

\tools Tome~1, Chapter~3 (binomial coefficients modulo $p$); Problem~\ref{pb:nullstellensatz}.
\leads Cauchy proved (a) in 1813 to fill a gap in Lagrange's four squares theorem, and Davenport rediscovered it in 1935. Vosper
(1956) showed that equality with $|A|,|B|\ge2$ and $|A+B|\le p-2$ forces arithmetic progressions with the same difference. Over
$\Z$ and other groups the subject became additive combinatorics.
\end{solution}

\begin{solution}{pb:erdosheil}
\idea Multiply the Cauchy--Davenport polynomial by $x-y$ to kill the diagonal $a=a'$, which is exactly what ``distinct'' removes.
The price is that the key coefficient is now a difference of two binomial coefficients, which is a Catalan number.

\writeup (a) Let $k=|A|$. For $k\le2$ the bound is trivial. If $2k-3\ge p$, then $p$ is odd (for $p=2$, $k\le2$) and $A$ contains a
subset $A'$ with $|A'|=\frac{p+3}2$, so it suffices to prove $|A'\dotplus A'|\ge p$; so we may assume $2k-3\le p$, hence
$2k-4<p$. Suppose $|A\dotplus A|\le2k-4$ and take $C\supset A\dotplus A$ with $|C|=2k-4$. Let
\[
f(x,y)=(x-y)\prod_{c\in C}(x+y-c).
\]
It vanishes on $A\times A$: on the diagonal because of $x-y$, off it because $x+y\in C$. Its degree is $2k-3=(k-1)+(k-2)$, and the
coefficient of $x^{k-1}y^{k-2}$ in $(x-y)(x+y)^{2k-4}$ is
\[
\binom{2k-4}{k-2}-\binom{2k-4}{k-1}=\binom{2k-4}{k-2}\Bigl(1-\frac{k-2}{k-1}\Bigr)=\frac1{k-1}\binom{2k-4}{k-2},
\]
the Catalan number $C_{k-2}$. Since $2k-4<p$, $p\nmid\binom{2k-4}{k-2}$, so $p\nmid C_{k-2}$. Problem~\ref{pb:nullstellensatz}(b) with
$S_1=S_2=A$, $t_1=k-1$, $t_2=k-2$, gives a point of $A\times A$ where $f\ne0$: a contradiction.

(b) For $A=\{0,1,\dots,k-1\}$ the sums of distinct elements are $1,2,\dots,2k-3$, which are $2k-3$ distinct residues when
$2k-3\le p$, that is $k\le\frac{p+3}2$.

\tools Problems~\ref{pb:nullstellensatz} and~\ref{pb:cauchydav}; Tome~0, Chapter~1 (binomial and Catalan numbers).
\leads Erdős and Heilbronn conjectured (a) in the 1960s; Dias da Silva and Hamidoune proved it in 1994 with exterior algebra, and
the polynomial proof above appeared a year later. The same polynomial with $\prod_{i<j}(x_i-x_j)$ bounds sums of $h$ distinct
elements: $|h^\wedge A|\ge\min(p,h|A|-h^2+1)$. The notebook checks (a) for all subsets when $p\le11$.
\end{solution}

\begin{solution}{pb:egz}
\idea Chevalley--Warning turns ``an equation of low degree has a non-trivial zero'' into a congruence. The polynomial $x^{p-1}$
is the indicator of $x\ne0$, so the two equations say ``choose $p$ indices whose $a_i$ add to $0$''.

\writeup (a) For $y\in\mathbb F_p$, $1-y^{p-1}$ is $1$ if $y=0$ and $0$ otherwise (Fermat). So the number $Z$ of common zeros
satisfies, in $\mathbb F_p$,
\[
Z\equiv\sum_{x\in\mathbb F_p^m}\prod_{j=1}^r\bigl(1-f_j(x)^{p-1}\bigr).
\]
The product is a polynomial of degree at most $(p-1)\sum_j\deg f_j<(p-1)m$, so each of its monomials $x_1^{e_1}\cdots x_m^{e_m}$ has
some $e_i<p-1$. Its sum over $\mathbb F_p^m$ is $\prod_i\sum_{x_i}x_i^{e_i}$, and $\sum_{x\in\mathbb F_p}x^e=0$ for $0\le e<p-1$
(Tome~2, Exercise~9.6, with $0^0=1$). So $Z\equiv0\pmod p$.

(b) Let $a_1,\dots,a_{2p-1}$ be the integers, read in $\mathbb F_p$. The system $f_1=f_2=0$ in $m=2p-1$ variables has degrees
$2(p-1)<m$, and $x=0$ is a common zero; by (a) there is another, $x\ne0$. Let $I=\{i:x_i\ne0\}$. Then $x_i^{p-1}=1$ for $i\in I$ and
$0$ otherwise, so $f_2(x)=|I|\equiv0$ and $f_1(x)=\sum_{i\in I}a_i\equiv0$. Since $0<|I|\le2p-1$, $|I|=p$.

(c) Take $2mn-1$ integers. As long as at least $2n-1$ remain, the theorem for $n$ lets us remove $n$ of them with sum divisible by
$n$. After $j$ removals $2mn-1-jn$ remain, which is at least $2n-1$ for $j\le2m-2$; so we get $2m-1$ disjoint groups, with sums
$nb_1,\dots,nb_{2m-1}$. The theorem for $m$ gives $m$ of the $b_i$ with sum divisible by $m$, and the union of these $m$ groups is a
set of $mn$ integers with sum divisible by $mn$. Since the theorem holds for $1$ and for primes, it holds for all $n$.

(d) Take $n-1$ zeros and $n-1$ ones: $n$ of them contain between $1$ and $n-1$ ones, so their sum is not divisible by $n$.

\tools Tome~2, Exercises~9.6 and~9.8 (power sums, Chevalley--Warning for one polynomial); Tome~1, Chapter~3 (Fermat).
\leads The original 1961 proof used Cauchy--Davenport instead. In the plane, Reiher proved in 2007 Kemnitz's conjecture: among
any $4n-3$ lattice points there are $n$ whose centroid is a lattice point.
\end{solution}

\begin{solution}{pb:hyperplanes}
\idea If fewer than $n(q-1)$ hyperplanes did the job, the product of their equations would agree, on all of $\mathbb F_q^n$, with
a polynomial whose leading monomial is $x_1^{q-1}\cdots x_n^{q-1}$, and the Nullstellensatz forbids this.

\writeup (a) The $n(q-1)$ hyperplanes $x_j=s$, $1\le j\le n$, $s\in\mathbb F_q^*$, cover every $x\ne0$ (some $x_j\ne0$) and miss $0$.
Suppose $m<n(q-1)$ hyperplanes $H_i=\{x:a_i\cdot x=b_i\}$ cover $\mathbb F_q^n\setminus\{0\}$ and miss $0$; then $b_i\ne0$ and we
may take $b_i=1$. Let
\[
f(x)=\prod_{i=1}^m(1-a_i\cdot x)-c\prod_{j=1}^n\prod_{s\in\mathbb F_q^*}(x_j-s),
\]
with $c$ chosen so that $f(0)=0$ (the second product is non-zero at $0$). For $x\ne0$, $x$ lies on some $H_i$, so the first product
vanishes, and some $x_j\in\mathbb F_q^*$, so the second does too. Thus $f$ vanishes on $\mathbb F_q^n$. But $\deg f=n(q-1)$, since
$m<n(q-1)$, and the coefficient of $x_1^{q-1}\cdots x_n^{q-1}$ is $-c\ne0$. Problem~\ref{pb:nullstellensatz}(b) with $S_j=\mathbb F_q$ and
$t_j=q-1$ gives a point where $f\ne0$: a contradiction.

(b) Let $B$ be a blocking set; translating, we may assume $0\in B$. Every hyperplane $\{x:a\cdot x=1\}$, $a\ne0$, misses $0$, so it
contains some $y\in B\setminus\{0\}$. In other words the sets $\{a:a\cdot y=1\}$, for $y\in B\setminus\{0\}$, are hyperplanes (in the
space of the $a$'s) missing $0$ and covering every $a\ne0$. By (a), $|B|-1\ge n(q-1)$. The union of the $n$ coordinate axes, with
$1+n(q-1)$ points, is blocking: a hyperplane $a\cdot x=b$ with $b\ne0$ meets the axis $j$ for any $a_j\ne0$, and one with $b=0$
contains $0$.

\tools Problem~\ref{pb:nullstellensatz}; Tome~0, Chapter~7 (hyperplanes).
\leads Over $\R$ the cube $\{0,1\}^n$ minus a point needs $n$ hyperplanes, and the same proof works. Projective blocking sets are
different: in a projective plane of order $q$ a line has $q+1$ points, and the real question is how small a blocking set
containing no line can be (Blokhuis 1994: at least $\frac32(q+1)$ for $q$ prime).
\end{solution}

\begin{solution}{pb:kakeya}
\idea Linear algebra gives a low-degree polynomial vanishing on a small $K$. On each line of $K$ it vanishes at $q$ points, more
than its degree, so it vanishes identically there, and its top coefficient along the line, which is the top homogeneous part
evaluated at the direction, is $0$. A polynomial of low degree that vanishes at every direction is $0$.

\writeup (a) The polynomials of degree at most $q-1$ in $n$ variables form a space of dimension $\binom{q-1+n}n$, the number of
monomials. If $|K|<\binom{q+n-1}n$, the linear map $f\mapsto(f(x))_{x\in K}$ has a non-trivial kernel: some $f\ne0$ of degree
$d\le q-1$ vanishes on $K$. Since $K\ne\emptyset$, $f$ is not a non-zero constant, so $d\ge1$. Let $f_d\ne0$ be its homogeneous part of
degree $d$. For $v\ne0$, take $a$ with $a+tv\in K$ for all $t$. Then $t\mapsto f(a+tv)$ is a polynomial in $t$ of degree at most $d<q$
with $q$ roots, so it is $0$, and its coefficient of $t^d$, which is $f_d(v)$, is $0$. Also $f_d(0)=0$. So $f_d$ vanishes on
$\mathbb F_q^n$ while $\deg_{x_i}f_d\le q-1<q$; by Problem~\ref{pb:nullstellensatz}(a), $f_d=0$. Contradiction. Finally
$\binom{q+n-1}n=\frac{q(q+1)\cdots(q+n-1)}{n!}\ge\frac{q^n}{n!}$.

(b) For $m\in\mathbb F_q$ let $L_m$ be the line $y=mx+m^2$, of slope $m$, and let $K$ be the union of the $L_m$ and the vertical line
$x=0$. A point $(x,y)$ lies on some $L_m$ when $m^2+mx-y=0$ has a root, that is when $x^2+4y$ is a square (zero included). For fixed
$x$, $y\mapsto x^2+4y$ is a bijection of $\mathbb F_q$, and $\frac{q+1}2$ elements are squares; so $\bigcup L_m$ has $\frac{q(q+1)}2$
points. On $x=0$ these are the points $(0,y)$ with $y$ a square, $\frac{q+1}2$ of them, so the vertical line adds $\frac{q-1}2$.

\tools Tome~0, Chapter~7 (dimension counting); Problem~\ref{pb:nullstellensatz}(a). Part~(a) is also Computing track~C1, Exercise~11.10; the construction of (b) is new here.
\leads Dvir's proof (2008) took one page and settled a problem that Wolff had posed in 1999 as a model for the real Kakeya
conjecture, still open in dimensions $\ge4$ (Wang and Zahl announced dimension $3$ in 2025). For $n=2$ and $q$ odd the
construction (b) is optimal (Blokhuis and Mazzocca, 2008); for $q=3$ the notebook confirms by brute force that the minimum is $7$,
one more than Dvir's bound $6$.
\end{solution}

\begin{solution}{pb:moore}
\idea The Frobenius $x\mapsto x^q$ is $\mathbb F_q$-linear. So the Moore matrix has dependent rows whenever the $x_i$ are
dependent over $\mathbb F_q$, and this gives all the linear factors of the determinant.

\writeup (a) Induction on $n$; for $n=1$, $\Delta=x_1$. Expanding along the last row, $\Delta(x_1,\dots,x_n)$ is a polynomial in
$x_n$ of degree $q^{n-1}$ with leading coefficient $\Delta_{n-1}=\Delta(x_1,\dots,x_{n-1})$. For $c\in\mathbb F_q^{n-1}$, substitute
$x_n=-\sum c_jx_j$: since $(\sum c_jx_j)^{q^k}=\sum c_jx_j^{q^k}$ (Frobenius is additive and fixes $\mathbb F_q$), the last row becomes
$-\sum c_j(\text{row }j)$, and the determinant vanishes. So each $x_n+\sum c_jx_j$ divides $\Delta$ in
$\mathbb F_q(x_1,\dots,x_{n-1})[x_n]$. These $q^{n-1}$ linear factors are pairwise distinct, so their product, of degree $q^{n-1}$ in
$x_n$, divides $\Delta$, and comparing leading coefficients
\[
\Delta(x_1,\dots,x_n)=\Delta_{n-1}\prod_{c\in\mathbb F_q^{n-1}}\Bigl(x_n+\sum_jc_jx_j\Bigr).
\]
The induction hypothesis on $\Delta_{n-1}$ finishes the proof (the identity holds in $\mathbb F_q[x_1,\dots,x_n]$ since both sides lie
there).

(b) Evaluate (a) at $\alpha$: $\Delta(\alpha)=0$ if and only if $\alpha_i+\sum_{j<i}c_j\alpha_j=0$ for some $i$ and some
$c\in\mathbb F_q^{i-1}$, that is, if and only if the $\alpha_i$ are linearly dependent over $\mathbb F_q$.

(c) Let $\alpha_1,\dots,\alpha_n$ be a basis of $V$. By (a),
\[
\Delta(\alpha_1,\dots,\alpha_n,X)=\Delta(\alpha)\prod_{v\in V}(X+v),
\]
and $\prod_{v\in V}(X+v)=\prod_{v\in V}(X-v)$ since $V=-V$. Expanding the determinant along its last row, $(X,X^q,\dots,X^{q^n})$,
gives a combination of $X,X^q,\dots,X^{q^n}$. Divide by $\Delta(\alpha)\ne0$. Each $X^{q^k}$ is additive, so $P_V$ is too.

\tools Tome~2, Chapter~9 (Frobenius); Tome~0, Chapter~7 (determinants).
\leads Moore's determinant is the $q$-analogue of the Vandermonde determinant. Linearised polynomials built this way give the
Gabidulin codes used in network coding, and Dickson's invariants of $GL_n(\mathbb F_q)$, the coefficients of
$\prod_{v\in\mathbb F_q^n}(X-v)$ written in the $x_i$.
\end{solution}

\begin{solution}{pb:capset}
\idea A diagonal tensor with $|A|$ non-zero entries has slice rank $|A|$, just as a diagonal matrix has rank $|A|$. A cap makes
$[x+y+z=0]$ diagonal on $A^3$. But the same function is a polynomial of degree $2n$, and in every monomial one variable has degree
at most $\frac{2n}3$: that writes it with few slices.

\writeup (a) Let $W$ have dimension $m$, and let $h\in W$ have the largest support $S$. If $|S|<m$, the subspace of $W$ of vectors
vanishing on $S$ is cut out by $|S|$ linear conditions, so it has dimension at least $m-|S|>0$ and contains some $w\ne0$. Since $w$
vanishes on $S$ and $h$ vanishes outside $S$, the support of $h+w$ is $S\cup\operatorname{supp}w\supsetneq S$, a contradiction.

(b) $T$ is the sum of the $|A|$ slices $c_a[x=a]\cdot[y=z=a]$, so its slice rank is at most $|A|$. Suppose
\[
T=\sum_{i=1}^{r_1}f_i(x)g_i(y,z)+\sum_{i=1}^{r_2}f'_i(y)g'_i(x,z)+\sum_{i=1}^{r_3}f''_i(z)g''_i(x,y).
\]
The functions $h:A\to F$ with $\sum_zf''_i(z)h(z)=0$ for all $i$ form a subspace of dimension at least $|A|-r_3$; by (a) one of
them has at least $|A|-r_3$ non-zero values. Multiply by $h(z)$ and sum over $z$. The left side becomes the diagonal matrix
$(x,y)\mapsto c_xh(x)[x=y]$, of rank at least $|A|-r_3$. The third sum disappears, and the first two become
$\sum_if_i(x)\tilde g_i(y)+\sum_i\tilde g'_i(x)f'_i(y)$, of rank at most $r_1+r_2$. So $r_1+r_2+r_3\ge|A|$.

(c) If $x+y+z=0$ and $x=y$, then $z=-2x=x$. If $x,y,z$ are pairwise distinct the cap condition forbids $x+y+z=0$. And if exactly
two coincide, say $x=y\ne z$, the first remark shows $x+y+z\ne0$. So on $A^3$, $[x+y+z=0]=[x=y=z]$, which is the $T$ of (b) with
$c_a=1$ (here $F=\mathbb F_3$).

(d) For $t\in\mathbb F_3$, $1-t^2$ is $1$ if $t=0$ and $0$ otherwise. Expanding the product gives a combination of monomials
$x^\alpha y^\beta z^\gamma$ with $\alpha,\beta,\gamma\in\{0,1,2\}^n$ and $|\alpha|+|\beta|+|\gamma|\le2n$, where
$|\alpha|=\sum\alpha_i$. In each monomial, one of $|\alpha|,|\beta|,|\gamma|$ is at most $\frac{2n}3$. Gather the monomials with
$|\alpha|\le\frac{2n}3$ by the value of $\alpha$: they form $\sum_\alpha x^\alpha G_\alpha(y,z)$, at most $M_n$ slices. Then those
remaining with $|\beta|\le\frac{2n}3$, then the rest, with $|\gamma|\le\frac{2n}3$. So the slice rank is at most $3M_n$, and by (c),
$|A|\le3M_n$.

(e) For $0<t\le1$, each $\alpha$ with $|\alpha|\le\frac{2n}3$ satisfies $t^{|\alpha|-2n/3}\ge1$, so
\[
M_n\le\sum_{\alpha\in\{0,1,2\}^n}t^{|\alpha|-2n/3}=\Bigl(\frac{1+t+t^2}{t^{2/3}}\Bigr)^n.
\]
The logarithmic derivative of $\frac{1+t+t^2}{t^{2/3}}$ vanishes when $\frac{1+2t}{1+t+t^2}=\frac2{3t}$, that is $4t^2+t-2=0$,
$t=\frac{\sqrt{33}-1}8=0.5931$, where the value is $\theta=2.7551$. So $|A|\le3\theta^n<3\cdot2.756^n$.

\tools Tome~0, Chapter~7 (rank, dimension); Tome~3, Chapter~4 (the exponential moment trick of Chernoff in (e)).
\leads The bound $\theta$ is optimal for this method, and the true growth rate is unknown: the best constructions give caps of size
about $2.22^n$ (Edel 2004; in 2023 a computer search guided by a language model, FunSearch, found a cap of $512$ points in $\mathbb F_3^8$).
The known maximum caps for $n=1,\dots,6$ are $2,4,9,20,45,112$, against $3M_n=3,9,30,45,153,504$. The slice rank method has since
proved the sunflower-free conjecture of Erdős and Szemerédi and bounded many other configurations.
\end{solution}

\begin{solution}{pb:paley}
\idea Adjacency is a value of the quadratic character, so every count in the Paley graph is a quadratic character sum, and those
are all $O(\sqrt p)$ by Problem~\ref{pb:conics}.

\writeup (a) The number of common neighbours of $0$ and $d\ne0$ is
\[
\sum_{z\ne0,d}\frac{1+\chi(z)}2\cdot\frac{1+\chi(z-d)}2=\frac14\Bigl[(p-2)+\sum_{z\ne0,d}\chi(z)+\sum_{z\ne0,d}\chi(z-d)+\sum_z\chi(z^2-dz)\Bigr].
\]
The two middle sums are $-\chi(d)$ and $-\chi(-d)=-\chi(d)$, and the last is $-1$ by Problem~\ref{pb:conics}(a) (discriminant $d^2\ne0$).
So it is $\frac{p-3-2\chi(d)}4$: $\frac{p-5}4$ if $d$ is a square (adjacent), $\frac{p-1}4$ otherwise. By translation the same holds for
any two vertices. Each of the $\frac{p(p-1)}4$ edges lies in $\frac{p-5}4$ triangles and each triangle has $3$ edges, so there are
$\frac{p(p-1)(p-5)}{48}$ triangles.

(b) For $s\ne t$, $[s\sim t]=\frac{1+\chi(s-t)}2$, and $\chi(0)=0$, so
$e(S,T)=\frac{|S||T|-|S\cap T|}2+\frac12\sum_{s\in S,t\in T}\chi(s-t)$. By Cauchy--Schwarz, extending the sum over $s$ to all of
$\mathbb F_p$,
\[
\Bigl|\sum_{s\in S}\sum_{t\in T}\chi(s-t)\Bigr|^2\le|S|\sum_{s\in\mathbb F_p}\Bigl|\sum_{t\in T}\chi(s-t)\Bigr|^2
=|S|\sum_{t,t'\in T}\sum_s\chi\bigl((s-t)(s-t')\bigr).
\]
The inner sum is $p-1$ if $t=t'$ and $-1$ otherwise (Problem~\ref{pb:conics}(a) again), so the total is
$|S|\bigl(|T|(p-1)-|T|(|T|-1)\bigr)=|S||T|(p-|T|)\le p|S||T|$.

(c) Let $K$ be a clique and $n$ a non-square. If $c+nc'=d+nd'$ with $c,c',d,d'\in K$ and $c\ne d$, then $c'\ne d'$, and
$c-d=n(d'-c')$ is a square times a non-square: a non-square, contradicting $c\sim d$. So $c=d$, then $c'=d'$: the map
$(c,c')\mapsto c+nc'$ is injective on $K\times K$, and $|K|^2\le p$. The map $x\mapsto nx$ is an isomorphism from $P_p$ to its
complement, so independent sets are bounded in the same way.

(d) $\sqrt{17}<5$, so by (c) $P_{17}$ has no clique and no independent set of size $5$. A direct check (done in the notebook)
shows that the largest ones have size $3$. A colouring of the edges of $K_{17}$ by ``square'' and
``non-square'' has no monochromatic $K_4$, so $R(4,4)>17$.

\tools Problem~\ref{pb:conics}; Tome~3, Chapter~2 (Cauchy--Schwarz); Tome~1, Chapter~6.
\leads Greenwood and Gleason (1955) proved $R(4,4)=18$, with $P_{17}$ as the extremal graph. The discrepancy bound (b) is the
definition of a \emph{quasirandom} graph in the sense of Chung, Graham and Wilson (1989): such graphs have about as many copies of
every small subgraph as a random graph. The clique bound $\sqrt p$ is far from the truth, which is believed to be $O(\log^2p)$,
the size of the cliques of a random graph being about $2\log_2p$; improving $\sqrt p$ even to $\sqrt{p/2}$ came only in 2021
(Hanson and Petridis).
\end{solution}

\begin{solution}{pb:mseq}
\idea $s_k$ is a fixed linear form evaluated at $\alpha^k$, and $\alpha^k$ runs over all of $\mathbb F_{2^n}^*$. Every statistic
of the sequence becomes a statistic of a linear form on the whole field, which is perfectly balanced.

\writeup (a) $\operatorname{Tr}$ is $\mathbb F_2$-linear (Frobenius is additive). From $f(\alpha)=0$, $\alpha^{k+n}=\sum_ic_i\alpha^{k+i}$,
and applying $\operatorname{Tr}$ gives the recurrence (signs do not matter in characteristic $2$).

(b) The trace is a polynomial of degree $2^{n-1}$, so it is not zero on all $2^n$ elements; it is a non-zero linear form onto
$\mathbb F_2$. If $\operatorname{Tr}(\beta\alpha^i)=0$ for $0\le i<n$, then $\operatorname{Tr}(\beta\gamma)=0$ for all $\gamma$, since the
$\alpha^i$ span $\mathbb F_{2^n}$ (it has degree $n$); if $\beta\ne0$, $\beta\gamma$ runs over the whole field, which is impossible.
So the linear map is injective, hence bijective. The window at $k$ is the image of $\beta=\alpha^k$. As $k$ runs over
$0,\dots,2^n-2$, $\alpha^k$ runs over the non-zero elements once each, so the windows are the non-zero $n$-bit blocks, once each.
Distinct windows mean the period is exactly $2^n-1$ (a window determines all later terms by (a), and $\alpha^{2^n-1}=1$). A sequence
satisfying (a) is determined by its first window; if that window is non-zero it equals the window of $(s_k)$ at some $k_0$, and then
the sequence is $(s_{k+k_0})$.

(c) The ones in a period are the non-zero $\beta$ with $\operatorname{Tr}\beta=1$; the kernel of $\operatorname{Tr}$ has $2^{n-1}$
elements, so $2^{n-1}$ have trace $1$, and $2^{n-1}-1$ non-zero ones have trace $0$. For the correlation,
$s_k+s_{k+t}=\operatorname{Tr}(\alpha^k(1+\alpha^t))$ with $1+\alpha^t\ne0$; as $k$ varies, $\alpha^k(1+\alpha^t)$ runs over the
non-zero elements, and
\[
\sum_{\beta\ne0}(-1)^{\operatorname{Tr}\beta}=\sum_\beta(-1)^{\operatorname{Tr}\beta}-1=0-1=-1.
\]

\tools Tome~2, Chapter~9 (the cyclic group $\mathbb F_q^*$, the trace; Exercise~9.3 for primitive polynomials); Tome~0, Chapter~7.
\leads These are Golomb's three randomness postulates (1967): balance, runs, two-valued autocorrelation. For $x^4+x+1$ the period is
$000100110101111$. The sequence is useless for cryptography, since $2n$ bits reveal the recurrence (Berlekamp--Massey), but its
flat correlation is what GPS satellites use: each one transmits a Gold code, the sum of two $m$-sequences of degree $10$.
\end{solution}

\begin{solution}{pb:singer}
\idea Multiplication by $g$ permutes the points of the projective plane $\mathbb P(\mathbb F_{q^3})$ cyclically. Two lines meet in
one point, and a line is never mapped to itself by a non-trivial power of $g$: that is the whole proof.

\writeup (a) $g^i$ and $g^{i+N}$ differ by the factor $g^N$, which has order $q-1$ and so generates $\mathbb F_q^*$; hence
$g^i\in U\iff g^{i+N}\in U$ and $D$ is well defined. The map $i\mapsto\mathbb F_qg^i$ is a bijection from $\Z/N\Z$ to the
$\frac{q^3-1}{q-1}=N$ points (lines through $0$), and $U$ contains $\frac{q^2-1}{q-1}=q+1$ of them: $|D|=q+1$.

Now fix $d\not\equiv0\pmod N$. The pairs $(i,j)\in D^2$ with $i-j=d$ correspond to the $j$ with $g^j\in U$ and $g^{j+d}\in U$, that is
to the points of $U\cap g^{-d}U$. Let $c=g^d$; then $c\notin\mathbb F_q$, because $\mathbb F_q^*=\langle g^N\rangle$. So
$\mathbb F_q(c)$ is a subfield of $\mathbb F_{q^3}$ strictly bigger than $\mathbb F_q$, and since its degree divides $3$ it is
$\mathbb F_{q^3}$. If $cU=U$, then $U$ would be stable under $\mathbb F_q[c]=\mathbb F_{q^3}$, a sub-$\mathbb F_{q^3}$-space of dimension
$2$ over $\mathbb F_q$, impossible. So $g^{-d}U$ and $U$ are distinct $2$-dimensional subspaces of a $3$-dimensional space; their
intersection has dimension $2+2-3=1$, a single point. So $d$ is a difference in exactly one way.

(b) Two points $x\ne y$ lie on $D+t$ when $x-t,y-t\in D$; their difference $x-y$ is a difference of $D$ in exactly one way, which
determines $x-t$, hence $t$. Two translates $D+t$ and $D+t'$, $t\ne t'$, meet at the points $x$ with $x-t$ and $x-t'$ in $D$; the
difference $t'-t$ is represented once, so there is one such $x$. (The translates are exactly the images $g^{-t}U$, the lines of the
plane.)

(c) Modulo $7$, the differences of $\{0,1,3\}$ are $\pm1,\pm2,\pm3$. Modulo $13$, those of $\{0,1,3,9\}$ are
\[
\pm1,\ \pm3,\ \pm9,\ \pm2,\ \pm8,\ \pm6,
\]
which are $\pm1,\dots,\pm6$ since $9\equiv-4$ and $8\equiv-5$. Modulo $21$, the differences of $\{0,1,4,14,16\}$ are
\[
\pm1,\ \pm4,\ \pm14,\ \pm16,\ \pm3,\ \pm13,\ \pm15,\ \pm10,\ \pm12,\ \pm2,
\]
and with $14\equiv-7$, $16\equiv-5$, $13\equiv-8$, $15\equiv-6$, $12\equiv-9$ these are $\pm1,\dots,\pm10$.

\tools Tome~2, Chapter~9 (subfields of $\mathbb F_{q^3}$, cyclicity of $\mathbb F_q^*$); Tome~0, Chapter~7 (dimension of an
intersection).
\leads Every known projective plane of order $n$ has $n$ a prime power, and every known cyclic one comes from Singer's construction;
the prime power conjecture for cyclic planes has been checked up to orders in the billions. The notebook rebuilds Singer sets for
$q=2,3,5,7$ from an explicit $\mathbb F_{q^3}$.
\end{solution}

\begin{solution}{pb:sidon}
\idea A Sidon set is a set with distinct differences. There are only $n-1$ possible differences, which gives $\sqrt{2n}$; counting
only short differences, in windows, saves the factor $\sqrt2$. Singer sets reach the bound.

\writeup (a) $a+b=a'+b'$ is the same as $a-a'=b'-b$, so in a Sidon set the $\binom k2$ positive differences are distinct, and they
lie in $\{1,\dots,n-1\}$. So $k(k-1)\le2(n-1)$ and $k<\sqrt{2n}+1$.

(b) Let $a_k=2pk+r_k$ with $r_k=k^2\bmod p\in[0,p)$. Then $a_k+a_l=2p(k+l)+(r_k+r_l)$ with $0\le r_k+r_l<2p$, so $a_k+a_l=a_{k'}+a_{l'}$
forces $k+l=k'+l'$ and $r_k+r_l=r_{k'}+r_{l'}$, hence $k^2+l^2\equiv k'^2+l'^2\pmod p$. With $s=k+l$, $2kl=s^2-(k^2+l^2)$, so
$kl\equiv k'l'$, and $\{k,l\}$, $\{k',l'\}$ are both the roots of $X^2-sX+kl$ in $\mathbb F_p$; as $0\le k,l,k',l'<p$, they are the
same pair. The largest element is below $2p(p-1)+p<2p^2$.

(c) Let $A\subset\{1,\dots,n\}$ be Sidon, $|A|=k$, and $1\le m\le n$. For $j=-m+1,\dots,n-1$ (that is $n+m-1$ values) let
$I_j=\{j+1,\dots,j+m\}$ and $A_j=|A\cap I_j|$. Each $a\in A$ lies in exactly $m$ windows, so $\sum_jA_j=km$, and by Cauchy--Schwarz
\[
\sum_j\binom{A_j}2\ge\frac12\Bigl(\frac{(km)^2}{n+m-1}-km\Bigr).
\]
On the other hand $\sum_j\binom{A_j}2$ counts pairs $a<b$ of $A$ with a common window, with multiplicity: if $b-a=d<m$ they share $m-d$
windows. Each $d$ is the difference of at most one pair, so $\sum_j\binom{A_j}2\le\sum_{d=1}^{m-1}(m-d)=\frac{m(m-1)}2$. Combining,
$\frac{k^2m}{n+m}-k\le m-1$, that is
\[
k^2\le\frac{(n+m)(m+k-1)}m=n\Bigl(1+\frac mn\Bigr)\Bigl(1+\frac{k-1}m\Bigr).
\]
Take $m=\lceil n^{3/4}\rceil$. By (a), $\frac{k-1}m\le\sqrt2\,n^{-1/4}$, so $k^2\le n\bigl(1+O(n^{-1/4})\bigr)$ and
$k\le\sqrt n+O(n^{1/4})$.

(d) Let $q$ be the largest prime with $q^2+q+1\le n$. The Singer set of Problem~\ref{pb:singer}, read as integers in $[0,N)$ and
shifted by $1$, is a Sidon set in $\{1,\dots,n\}$ with $q+1$ elements: distinct differences modulo $N$ are distinct differences
in $\Z$. By the prime number theorem (Tome~5, Chapter~5), consecutive primes satisfy $p_{j+1}/p_j\to1$, so $q\sim\sqrt n$. With
(c), $F(n)\sim\sqrt n$.

\tools Tome~3, Chapter~2 (Cauchy--Schwarz); Problem~\ref{pb:singer}; Tome~5 (the prime number theorem).
\leads Erdős offered \$500 for a proof or disproof of $F(n)=\sqrt n+O(1)$. The best upper bound is
$\sqrt n+0.998n^{1/4}$ for large $n$ (Balogh, Füredi and Roy, 2023), improving the constant $1$ of Lindström (1969). Sidon sets are
also the rulers with all distances distinct, \emph{Golomb rulers}, used to place radio antennas.
\end{solution}

\begin{solution}{pb:randpoly}
\idea Polynomials over $\mathbb F_q$ factor uniquely, so their counting functions satisfy Euler products, exactly like the
integers, but with power series in $t=q^{-s}$ that can be summed in closed form.

\writeup (a) Every monic $f$ is uniquely $g^2h$ with $g$, $h$ monic and $h$ squarefree (group the irreducible factors by parity of
their multiplicity). With $s_n$ the number of squarefree monic polynomials of degree $n$ and $S(t)=\sum s_nt^n$, this says
$\frac1{1-qt}=S(t)\cdot\frac1{1-qt^2}$, so $S(t)=\frac{1-qt^2}{1-qt}$ and $s_n=q^n-q^{n-1}$ for $n\ge2$.

(b) By Tome~2, Theorem~9.10, $nM(n,q)=\sum_{d\mid n}\mu(d)q^{n/d}$. The term $d=1$ is $q^n$, and the others have $n/d\le n/2$, so
$|nM(n,q)-q^n|\le\sum_{e\le n/2}q^e\le\frac q{q-1}q^{n/2}\le2q^{n/2}$. Divide by $nq^n$.

(c) Choosing $f$ of type $\lambda$ means choosing, for each $j$, a multiset of $m_j$ monic irreducibles of degree $j$: there are
$\prod_j\binom{M(j,q)+m_j-1}{m_j}$ such $f$. As $q\to\infty$, $M(j,q)\sim\frac{q^j}j$ by (b), so
$\binom{M+m-1}m\sim\frac{M^m}{m!}\sim\frac{q^{jm}}{j^mm!}$, and the product is $\sim q^{\sum jm_j}/\prod_jj^{m_j}m_j!=q^n/z_\lambda$. And
$1/z_\lambda$ is the probability that a uniform permutation of $\mathfrak S_n$ has cycle type $\lambda$, because the class of type
$\lambda$ has $n!/z_\lambda$ elements (Tome~2, Chapter~4).

\tools Tome~2, Chapter~9 (counting irreducibles); Tome~4, Chapter~2 (generating functions); Tome~2, Chapter~4 (conjugacy classes of
$\mathfrak S_n$).
\leads Compare with the integers: squarefree density $1/\zeta(2)=6/\pi^2$, and here $1/Z(q^{-2})=1-\frac1q$ where
$Z(t)=\frac1{1-qt}$ plays the role of $\zeta$. The error $q^{-n/2}$ in (b) is the Riemann hypothesis for $\mathbb F_q[t]$, which is
easy; for the integers it is the open problem. The analogy ``factor degrees = cycle lengths'' is the Frobenius acting on the roots,
and Chebotarev's theorem makes it precise.
\end{solution}

\begin{solution}{pb:noroots}
\idea Evaluation at the $q$ points of $\mathbb F_q$ is a linear map, onto as soon as the degree allows interpolation. The roots then
behave like $q$ independent coins of bias $\frac1q$, exactly when $n\ge q$, and like the first terms of inclusion--exclusion
otherwise.

\writeup Write $f=x^n+g$ with $g$ uniform among polynomials of degree $<n$. For a set $E\subset\mathbb F_q$ with $|E|\le n$, the linear map
$g\mapsto(g(a))_{a\in E}$ is onto (Lagrange interpolation with polynomials of degree $<|E|$), so it takes every value equally often,
and so does $f\mapsto(f(a))_{a\in E}$, which is a translate of it.

(a) If $n\ge q$, take $E=\mathbb F_q$: $(f(a))_a$ is uniform on $\mathbb F_q^q$, the coordinates are independent and uniform, the events
$f(a)=0$ are independent with probability $\frac1q$, and $R\sim\mathrm{Bin}(q,\frac1q)$. Hence $P(R=0)=(1-\frac1q)^q$.

(b) By inclusion--exclusion,
\[
P(R=0)=\sum_{E\subset\mathbb F_q}(-1)^{|E|}P(f=0\text{ on }E).
\]
For $|E|\le n$ this probability is $q^{-|E|}$ by the
remark above; for $|E|>n$ it is $0$, a monic polynomial of degree $n$ having at most $n$ roots. So
$P(R=0)=\sum_{k\le n}(-1)^k\binom qkq^{-k}$, and since
\[
\binom qkq^{-k}=\frac1{k!}\prod_{i<k}\Bigl(1-\frac iq\Bigr)\longrightarrow\frac1{k!},
\]
as $q\to\infty$ this tends to $\sum_{k\le n}\frac{(-1)^k}{k!}$, the derangement probability: a random polynomial has no root as often as a random
permutation has no fixed point, in agreement with Problem~\ref{pb:randpoly}(c), roots being factors of degree $1$.

\tools Tome~0, Chapter~3 (Lagrange interpolation); Tome~3, Chapter~1 (inclusion--exclusion, binomial law); Chapter~4 of this book
(derangements).
\leads For $n\ge q$ the answer does not depend on $n$ at all. Over $\Z$ the analogue is the number of roots modulo $p$ of a random
polynomial, and for a fixed polynomial the Chebotarev density theorem gives the law of the number of roots modulo a random prime.
\end{solution}

\begin{solution}{pb:waring}
\idea The number of $k$-th roots of $u$ is a sum over the $k$ characters of order dividing $k$. Then $N(c)$ is $p$ plus a sum of
$(k-1)^2$ Jacobi sums, of which $k-1$ have size $1$ and the others size $\sqrt p$.

\writeup (a) $\mathbb F_p^*$ is cyclic and $k\mid p-1$, so the characters $\chi$ with $\chi^k=\varepsilon$ form a group of order $k$,
and for $u\ne0$, $\sum_{\chi^k=\varepsilon}\chi(u)$ is $k$ if $u$ is a $k$-th power and $0$ otherwise: the number of $x$ with $x^k=u$.
With $\varepsilon(0)=1$ and $\chi(0)=0$ for $\chi\ne\varepsilon$ it is also right for $u=0$. So
\[
N(c)=\sum_{\chi,\psi}\sum_{u+v=c}\chi(u)\psi(v)=\sum_{\chi,\psi}\chi\psi(c)J(\chi,\psi),
\]
after the substitution $u=cu'$, $v=cv'$. The pair $(\varepsilon,\varepsilon)$ gives $p$; pairs with exactly one trivial character
give $0$; the $k-1$ pairs $(\chi,\bar\chi)$ with $\chi\ne\varepsilon$ give $J(\chi,\bar\chi)=-\chi(-1)$, of absolute value $1$ (as in
Problem~\ref{pb:fermatcubic}(a)); and the remaining $(k-1)^2-(k-1)=(k-1)(k-2)$ pairs have $\chi,\psi,\chi\psi$ all non-trivial,
so $|J(\chi,\psi)|=|g(\chi)||g(\psi)|/|g(\chi\psi)|=\sqrt p$ by Problem~\ref{pb:fermatcubic}(b), whose proof works for any characters.

(b) $c=0$ is $0^k+0^k$. For $c\ne0$, $N(c)>0$ if $p>(k^2-3k+2)\sqrt p+k-1$. If $p\ge k^4$, then $\sqrt p\ge k^2$ and
$p-(k^2-3k+2)\sqrt p-(k-1)\ge(3k-2)\sqrt p-(k-1)>0$.

(c) If $d=\gcd(k,p-1)$, the $k$-th powers of $\mathbb F_p^*$ are the $d$-th powers (the image of $x\mapsto x^k$ in the cyclic group
of order $p-1$ is the subgroup of index $d$). So the question reduces to $d$; when $d=1$, every residue is a $k$-th power.

\tools Tome~1, Chapter~5 (characters of $\mathbb F_p^*$); Problem~\ref{pb:fermatcubic}.
\leads For $k=3$ the bound is $2\sqrt p+2$, matching Problem~\ref{pb:fermatcubic}. Weil's theorem bounds the error by
$(k-1)(k-2)\sqrt p$ for the projective curve, of genus $\frac{(k-1)(k-2)}2$. The notebook lists the primes where two $k$-th powers
miss a residue: for $k=3$ only $p=7$, for $k=4$ the largest is $29$, for $k=8$ it is $761$. They are all well below $k^4$, so the
exponent $4$ is far from sharp.
\end{solution}

\section*{Map of the chapter}

\begin{center}
\small
\adjustbox{max width=\linewidth}{\begin{tabular}{lll}
Problem & Topic & Tomes used\\\hline
\ref{pb:conics} & points on conics & 1\\
\ref{pb:sumsquares} & sums of squares, Gauss sums & 1, 3\\
\ref{pb:fermatcubic} & Gauss's cubic, Jacobi sums & 1, 2\\
\ref{pb:nilpotent} & nilpotent matrices (Fine--Herstein) & 0, 2\\
\ref{pb:rankdist} & rank of a random matrix, Durfee & C1, 3, 4\\
\ref{pb:flags} & flags, the field with one element & 0, 2, C1\\
\ref{pb:schwartz} & Schwartz--Zippel, matchings & 0\\
\ref{pb:nullstellensatz} & Combinatorial Nullstellensatz & 0\\
\ref{pb:cauchydav} & Cauchy--Davenport & 1\\
\ref{pb:erdosheil} & Erdős--Heilbronn & 0\\
\ref{pb:egz} & Erdős--Ginzburg--Ziv & 1, 2\\
\ref{pb:hyperplanes} & covering by hyperplanes, blocking sets & 0\\
\ref{pb:kakeya} & Dvir's Kakeya theorem & 0\\
\ref{pb:moore} & Moore determinant, linearised polynomials & 0, 2\\
\ref{pb:capset} & the cap set theorem & 0, 3\\
\ref{pb:paley} & Paley graphs & 1, 3\\
\ref{pb:mseq} & maximal length sequences & 0, 2\\
\ref{pb:singer} & Singer difference sets & 0, 2\\
\ref{pb:sidon} & Sidon sets & 3, 5\\
\ref{pb:randpoly} & random polynomials and permutations & 2, 4\\
\ref{pb:noroots} & roots of a random polynomial & 0, 3\\
\ref{pb:waring} & two $k$-th powers & 1\\
\end{tabular}}
\end{center}

\section*{Sources and further reading}

\begin{itemize}
\item G.~H.~Hardy, \emph{A Mathematician's Apology}, Cambridge University Press, 1940, \S10 (the epigraph).
\item B.~L.~Davis and D.~Maclagan, The card game SET, \emph{The Mathematical Intelligencer} 25:3 (2003), 33--40 (Falco's
invention of the game, and Pellegrino's $20$-point cap of 1971).
\item R.~Meshulam, On subsets of finite abelian groups with no 3-term arithmetic progressions, \emph{Journal of Combinatorial
Theory A} 71 (1995), 168--172; M.~Bateman and N.~H.~Katz, New bounds on cap sets, \emph{Journal of the American Mathematical
Society} 25 (2012), 585--613.
\item K.~Ireland and M.~Rosen, \emph{A Classical Introduction to Modern Number Theory}, 2nd ed., Springer, 1990, Chapters~8 and~10
(Gauss and Jacobi sums, the equation $x^3+y^3=1$).
\item A.~Weil, Numbers of solutions of equations in finite fields, \emph{Bulletin of the American Mathematical Society} 55 (1949),
497--508.
\item R.~Lidl and H.~Niederreiter, \emph{Finite Fields}, Cambridge University Press, 1997 (linearised polynomials, the Moore
determinant, sequences).
\item N.~J.~Fine and I.~N.~Herstein, The probability that a matrix be nilpotent, \emph{Illinois Journal of Mathematics} 2 (1958),
499--504.
\item J.~Fulman, Random matrix theory over finite fields, \emph{Bulletin of the American Mathematical Society} 39 (2002), 51--85.
\item N.~Alon, Combinatorial Nullstellensatz, \emph{Combinatorics, Probability and Computing} 8 (1999), 7--29.
\item N.~Alon, M.~B.~Nathanson and I.~Ruzsa, The polynomial method and restricted sums of congruence classes, \emph{Journal of Number
Theory} 56 (1996), 404--417.
\item N.~Alon and Z.~Füredi, Covering the cube by affine hyperplanes, \emph{European Journal of Combinatorics} 14 (1993), 79--83.
\item Z.~Dvir, On the size of Kakeya sets in finite fields, \emph{Journal of the American Mathematical Society} 22 (2009), 1093--1097.
\item L.~Guth, \emph{Polynomial Methods in Combinatorics}, American Mathematical Society, 2016.
\item E.~Croot, V.~Lev and P.~P.~Pach, Progression-free sets in $\Z_4^n$ are exponentially small, \emph{Annals of Mathematics} 185
(2017), 331--337.
\item J.~S.~Ellenberg and D.~Gijswijt, On large subsets of $\mathbb F_q^n$ with no three-term arithmetic progression, \emph{Annals of
Mathematics} 185 (2017), 339--343.
\item T.~Tao, A symmetric formulation of the Croot--Lev--Pach--Ellenberg--Gijswijt capset bound, blog post, 18 May 2016.
\item B.~Romera-Paredes et al., Mathematical discoveries from program search with large language models, \emph{Nature} 625 (2024),
468--475.
\item F.~R.~K.~Chung, R.~L.~Graham and R.~M.~Wilson, Quasi-random graphs, \emph{Combinatorica} 9 (1989), 345--362.
\item S.~W.~Golomb, \emph{Shift Register Sequences}, Holden-Day, 1967.
\item J.~Singer, A theorem in finite projective geometry and some applications to number theory, \emph{Transactions of the American
Mathematical Society} 43 (1938), 377--385.
\item P.~Erdős and P.~Turán, On a problem of Sidon in additive number theory, and on some related problems, \emph{Journal of the
London Mathematical Society} 16 (1941), 212--215.
\item K.~O'Bryant, A complete annotated bibliography of work related to Sidon sequences, \emph{Electronic Journal of Combinatorics}
DS11 (2004).
\item A.~Granville, Cycle lengths in a permutation are typically Poisson, \emph{Electronic Journal of Combinatorics} 13 (2006), R107
(the analogy between polynomials and permutations).
\end{itemize}
